% 本文件由 LaTeX 排版 Agent 根据有效 translation.json 自动生成。
% author=Councils; work_id=3819; title=尼西亚第二会议

\chapter{尼西亚第二会议}

% structure: p0001 source=h1 target=omitted reason=page-title-already-rendered-by-parent

% structure: p0002 source=h2 target=section reason=relative-heading-rank; base=h2

\section{君士坦丁和依勒纳皇帝致最神圣、最蒙福的旧罗马教宗哈德良的神圣谕书}

（\Emph{收录于拉贝和科萨特}，\WorkTitle{Concilia}，\Emph{第七卷，第32栏}。）\par

凡从我们的主耶稣基督接受帝国尊位或首席司铎的尊荣的人，理应提供并关切那些中悦于祂的事，并按照祂的旨意和喜好治理管理托付给他们的人民。\par

因此，最神圣的\OriginalTerm{元首}{Caput}，我们和您都有责任，要无可指责地明白属于祂的事，并在此事上操练自己，因为我们从祂接受了帝国尊位，您则接受了首席司铎的尊荣。\par

但现在更切题地说吧。您的父辈之福知道，在此前我们这皇城中对可敬圣像的所作所为，我们之前统治的那些人如何毁坏它们，使它们蒙受耻辱和伤害（但愿这不算在他们的账上，因为他们若没有伸手攻击教会倒还好！），以及他们如何引诱并说服所有住在这片地区的人接受他们的意见——是的，整个东方也是如此，直到天主将我们提升到这个国度，我们真诚寻求祂的荣耀，并持守由祂的宗徒及其他所有导师所传下来的道理。因此，现在我们怀着纯洁的心和真诚的信仰，与所有臣民及最博学的神学家们，就有关天主的事进行了不断的商议，并根据他们的建议决定召开一次大公会议。我们恳请您的父辈之福，更确切地说是主天主恳请——\BibleQuote{祂愿意所有的人都得救，并得以认识真理}\BibleRef{弟前 2:4}——请您将自己赐予我们，不要迟延，而要上到这里来，帮助我们确认和确立可敬圣像的古老传统。确实，您的圣德应当这样做，因为您知道经上写着——\BibleQuote{我的百姓，你们司祭，要安慰，安慰他们，说：天主这样说}\BibleRef{依 40:1}，以及 \BibleQuote{司祭的嘴唇应守护知识，人也从他的口中寻求训诲，因为他是万军上主的天使}\BibleRef{拉 2:7}。还有，那位宣讲真理的神圣宗徒，\BibleQuote{从耶路撒冷及四周直到依里利孔，传遍了福音}\BibleRef{罗 15:19}，如此命令说：\BibleQuote{你们要以纪律牧养基督的羊群，祂以自己的血购得了他们}\BibleRef{宗 20:28}。既然您是真真实实的\OriginalTerm{首席司铎}{primus sacerdos}，主持在圣且极可赞颂的宗徒伯多禄的位置和教座上，愿您的父辈之福来我们这里，如同我们之前所说，并与其他所有将在此聚集的司铎一同出席，这样主的旨意就得以成就。因为我们在福音中得知主说——\BibleQuote{哪里有两个或三个人因我的名字聚在一起，我就在他们中间}\BibleRef{玛 18:20}——愿您的父辈和神圣之福被伟大的天主、万有的君王、我们的主耶稣基督以及我们这些祂的仆人所确证和保证：若您上到这里来，您将受到一切荣誉和荣耀的接待，所需的一切都将赐予您。再者，当\OriginalTerm{定义}{capitulum}完成时（我们希望在基督我们天主的恩许下完成），我们负责为您提供一切方便，使您得以光荣而卓越地返回。然而，若您的福分不能莅临我们这里（这我们几乎难以想象，因我们深知您对神圣事物的热忱），那么至少请为我们选派有智慧的人，由他们携带您圣德的信函，如此他们可以代表您神圣的父辈之福出席。如此，当他们与这里的其他司铎会面时，我们圣教父们的古老传统可藉以在会议中得到确认，所有邪恶的莠子都可被拔除，我们的主和救主耶稣基督的话便得以成就：\BibleQuote{地狱的权势不得战胜她}\BibleRef{玛 16:18}。此后，在唯一、神圣、至公、从宗徒传下来的教会中，再无分裂和分离，因为我们的真天主基督是她的头。\par

我们已将我们亲爱的西西里的莱昂蒂纳最神圣的主教、在主内蒙爱的君士坦丁（他与您的父辈之福很熟悉）召到我们面前，与他当面交谈后，随同我们这份可敬的谕旨派遣他到您那里。他在觐见您之后，请立即打发他回来，好让他回到我们这里，并让他告诉我们您到来的消息——从那里启程到我们这里需要多少时间。此外，他还可以与最神圣的那不勒斯主教一同留住，并与之一同上到这里。由于您的行程将经过那不勒斯和西西里，我们已就此命令西西里总督，要他妥善安排一切为您尊荣和休息所必需的准备，以便您的父辈之福能够来到我们这里。于\OriginalTerm{九月朔日前第四日}{ivth before the calends of September}（即八月二十九日），第七小纪，发自皇城。\par

% structure: p0008 source=h2 target=section reason=relative-heading-rank; base=h2

\section{帝国谕书}

\Emph{在第一会议宣读。}\par

（\Emph{收录于拉贝和科萨特}，\WorkTitle{Concilia}，\Emph{第七卷，第49栏}。）\par

君士坦丁和依勒纳——秉持信仰的罗马君主，致因天主的恩宠及我们虔诚治权的命令，聚集在尼西亚会议的最神圣主教们。\par

真正按照天主圣父本性的智慧——我们的主耶稣基督，我们的真天主——借着祂最神圣、最奇妙的肉身之安排，将我们从一切偶像崇拜的谬误中拯救出来；并借着承担我们的本性，借着与祂同本性的圣神的合作，使之更新；祂自己成为首位大司祭，并使你们圣人们堪当同等的尊荣。\par

祂就是那位善牧，亲自用祂的肩膀将那只迷失的羊——堕落的人类——背回自己特有的羊圈，即天使与服役之灵的团体\BibleRef{弗 2:14}，并将我们与祂自己合好，拆毁了中间隔断的墙壁，藉着祂的肉体消除了仇恨，并赐给我们和平的生活规范；因此，祂在福音中向众人宣讲说：\BibleQuote{缔造和平的人是有福的，因为他们要被称为天主的子女。}\BibleRef{玛 5:9}。我们敬虔的君主渴望有分于这祝福，这祝福确认了义子尊位的崇高，因此始终竭尽全力引导我们整个罗马帝国走向合一与和谐；我们尤其关心天主的教会的正确秩序，并千方百计促进司祭职的合一。为此，东方、北方、西方与南方的司祭职首领们，都藉着他们的代表主教出席，这些代表各自携带了回复至圣宗主教所发出的公函而写的答复；因为自起初，这便是在整个领受福音的天主教会内的教务会议规定。为此，凭天主的善意与允许，我们召集了你们，祂至圣的司祭们，你们惯于在无血之祭中分施祂的证言，为使你们的决定符合先前那些正确颁布法令的会议的界定，并使圣神的光辉在一切事上光照你们，因为正如我们的主教导的：\BibleQuote{没有人点灯放在斗底下，而是放在灯台上，照耀屋里所有的人。}同样，你们应善用我们祖先虔诚流传下来的各种规条，使天主的众圣教会都安享和平秩序。\par

至于我们，我们对真理的热忱如此之大——对宗教利益的渴望如此之切——对教会秩序的关怀如此之深——对古代规则与秩序得以维持的焦虑如此之重——尽管我们全神贯注于军事会议——尽管我们的注意力完全被政治事务占据——然而，我们视这些事务为次要，不允许任何事情妨碍你们至圣会议的召集。每个人都被赋予完全的自由，毫无迟疑地表达自己的意见，以便研究的主题得到充分讨论，真理得以更勇敢地言说，从而使一切纷争从教会中消除，我们都在和平的纽带中合而为一。\par

因为，当至圣宗主教\Person{保禄}，按天主的旨意，即将从死亡的束缚中解放，并离开尘世的旅途，与他的主\Person{基督}一同回天上的家乡时，他退下了宗主教职位，接受了隐修生活。我们问他：\Quote{你为什么这样做？}他回答说：\Quote{因为我害怕，如果死亡在这座皇家的、有天主保佑的城市的主教职位上突然捉住我，我将不得不带着整个天主教会的\OriginalTerm{绝罚}{anathema}同去，这绝罚将我投入那为魔鬼及其使者预备的外面的黑暗中；因为有人说：\InnerQuote{这里曾举行过一次会议，旨在推翻天主教会为纪念人物而持有、拥抱并接纳的图像。}这就是使我灵魂纷扰的原因——这就是使我焦虑地寻求如何逃脱天主的审判的原因——因为我是在这些人中长大的，并且被列于他们之中。}他刚在我们一些最显赫的贵族面前说完这些话，便去世了。\par

当我们虔诚的君主默思这可怕的宣告时（实在地，甚至在此之前，我们就已从许多周围的人那里听到类似的问题），我们与自己商议应当如何行；经过深思熟虑后，我们决定，待新宗主教当选后，应努力使此议题得到决定性结论。因此，我们召集了那些我们所知在教会事务上最有经验的人，并呼求我们的天主基督，与他们商议谁堪当被提升至这皇家的天主保佑之城的司祭职宝座；他们众口一心，一致投票赞成\Person{塔拉西乌斯}——即现任宗座首席者。于是我们派人召他来，将我们的商议和投票结果放在他面前；但他绝不同意，也丝毫不屈服于所决定的事。当我们询问他为何拒绝同意时，他起初回避回答说，首席司祭职的重担对他而言过于沉重。但我们知道这不过是他不愿服从我们的托词，并不放弃我们的坚持，继续迫使他接受首席司祭职的尊严。当他发现我们如此迫切时，便告诉我们他拒绝的原因。他说：\Quote{因为我看到那建立在磐石、我们的天主基督之上的教会，因分裂而撕裂破碎，我们在宣认上不稳定，东方的基督徒虽与我们有相同信仰，却拒绝与我们共融，而与西方之人联合；因此我们与所有人疏远，且每日被所有人诅咒；此外，我应要求召开一次大公会议，应有罗马教宗和东方首席司祭的使节出席。}因此，我们充分明白这些事后，将他引荐给聚集的司祭团——我们最杰出的王子们——以及全体基督徒百姓；然后，当着他们的面，他将先前对我们所说的一切又向他们重述；他们听后，欢喜地接纳了他，并热切地恳求我们爱好和平的虔诚君主召开大公会议。我们衷心同意他们的请求；因为说实话，是天主的美善和引导使我们将你们聚集在一起。因此，既然天主要成就祂自己的计划，为此目的将你们从世界各地聚集于此，看，福音书现在摆在你们面前，清楚大声地呼喊：\Quote{秉公审判；}你们要坚定地作宗教的捍卫者，随时准备以毫不留情的手砍除一切革新和标新立异的发明。正如宗徒团的首领\Person{伯多禄}以剑击打疯狂的奴仆，割下他的犹太耳朵一样，同样，你们要挥舞圣神之斧，每棵结出争论、纷争或新输入革新的树，要么通过健康教义的话语移植而更新，要么以教规的谴责将其砍倒，投入未来地狱之火，以使圣神的和平永远保护整个教会之体，使之巩固合一，并由圣传所确认；如此，我们整个罗马国家也与教会一样享受和平。\par

我们已收到最圣洁的旧罗马教宗\Person{哈德良}通过其使节——即天主所爱的总司铎\Person{伯多禄}和天主所爱的司铎兼院长\Person{伯多禄}——送来的信函，他们将与你们一同出席公会议；我们命令，按照主教会议惯例，这些信函应在你们全体面前宣读；并且，你们在适度的静默中聆听了这些，以及东方教区首席司祭和其他司祭通过最虔诚的隐修士兼安提约基雅宗主教宝座秘书长\Person{若望}和司铎兼院长\Person{多默}（他们也与你们一起出席）送来的两本八开本中所含的信函后，可以由此了解天主教会对此点的看法。\par

% structure: p0018 source=h2 target=section reason=relative-heading-rank; base=h2

\section{会议录摘录}

% structure: p0019 source=h3 target=subsection reason=relative-heading-rank; base=h2

\subsection{第一次会议}

（\WorkTitle{会议集}，第七卷，第53栏）\par

[\Emph{某些曾被圣像破坏者引入歧途的主教前来请求重新接纳。其中第一位是安西拉的\Person{巴西略}。}]\par

安西拉主教\Person{巴西略}从一本书中宣读如下：鉴于教会立法自古代以来，甚至从起初由圣宗徒及其继承人（即我们的圣父和导师）以及六次神圣大公会议和为维护正统信仰而召开的地方会议，均按照教规传下来，即凡从任何异端回归正统信仰和天主教会传统者，应否认自己的异端，并宣认正统信仰。\par

因此，我，安西拉城主教\Person{巴西略}，提议与天主教会、旧罗马最圣洁的教宗\Person{哈德良}、最蒙福的宗主教\Person{塔拉西乌斯}以及至圣的宗座，即亚历山大里亚、安提约基雅和圣城，以及所有正统大司祭和司祭合一，特写下此信仰宣认，并奉献给你们，如同献给那些因宗徒权威而领受权柄的人。在此，我也恳请你们天主所聚集的圣洁者，宽恕我在此事上的迟缓。因为我不应在正统宣认上落后，但这完全出于我对此事的一无所知，以及怠慢和疏忽的心思。因此，我更加恳请你们的蒙福，在天主面前赐予我宽容。\par

因此，我信并宣认在一天主、全能圣父，并在独一主耶稣基督、祂的唯一圣子，并在圣神、主及赋予生命者之内。天主圣三，一体同尊，应在一个天主性、德能与权威中被钦崇和荣耀。我宣认关于天主圣三之一位、我们的主天主耶稣基督降生成人的一切，正如圣人和六次大公会议所传授的。我弃绝并诅咒一切异端的妄言，正如他们也已弃绝。我恳求我们无玷之母、圣天主之母的\OriginalTerm{转求}{\GreekText{πρεσβείας}}，以及神圣天朝诸军的转求，和全体圣人的转求。\par

并且，我以一切\OriginalTerm{尊敬}{\GreekText{τιμῆς}}接受他们神圣可敬的遗骸，以\OriginalTerm{以尊敬之情敬礼}{\GreekText{τιμητικῶς} \GreekText{προσκυνέω}}向这些致敬并敬礼，盼望在他们的圣德中有份。同样，也可敬的\OriginalTerm{圣像}{\GreekText{εἰκόνας}}，即我们的主耶稣基督为我们的救恩所取的人性的圣像，以及我们无玷之母、圣天主之母的圣像，和相似天主的天使的圣像，以及圣宗徒、先知、殉道者、全体圣人的圣像——所有这些神圣的圣像，我致敬并敬礼——以我全部的灵魂和心思弃绝并诅咒那因顽梗和狂妄聚集起来的会议，它自称为第七次大公会议，但被正确思考的人合法且依规地称为伪会议，因为它违反一切真理和虔诚，鲁莽轻率地反对天主传授的教会法规，甚至不敬地狂吠并嘲讽神圣可敬的圣像，并下令将这些圣像从天主的圣教会中移除；主持那会议的是假称为以弗所人的\Person{德奥多西}、佩尔格的\Person{西西尼乌斯}（别号帕斯提拉斯）、彼西底亚的\Person{巴西略}（伪称\Quote{tricaccabus}）；悲惨的\Person{君士坦丁}（当时之宗主教）与他们一同\OriginalTerm{被迷惑}{\GreekText{ἐματαιώθη}}。\par

这些事我如此宣认并同意，因此，在天主面前，以心里的纯朴和正直的心志，我签署了下列绝罚条文。\par

弃绝诽谤基督徒者，即破坏圣像者。\par

弃绝将圣经中针对偶像的话语应用于可敬圣像者。\par

弃绝不向神圣可敬圣像致敬者。\par

弃绝说基督徒求助于圣像如同求助于神者。\par

弃绝称神圣圣像为偶像者。\par

弃绝明知而与诋毁、羞辱可敬圣像者共融者。\par

弃绝说另有他人而非我们的主基督将我们从偶像中救出者。\par

弃绝藐视圣教父的教导和天主教会的传统，以阿里乌、聂斯多略、欧迪克、狄奥斯哥鲁的论证为借口并自以为是者，声称除非我们显然受旧约和新约的教导，我们不应遵行圣教父和神圣大公会议的教导及天主教会的传统。\par

弃绝胆敢说天主教会曾许可偶像者。\par

弃绝说制造圣像是魔鬼的发明而非我们圣教父的传统者。\par

这是我的〔信仰〕宣认，我同意这些主张。我以全心、全灵、全意宣告此事。\par

若在任何时候，因魔鬼的欺诈（愿天主禁止！），我自愿或不自愿地反对我现在所宣认的，愿我受圣父、圣子和圣神的弃绝，并成为天主教会和一切圣统秩序的局外人。\par

我将按圣宗徒和公认教父的神圣教规，远离一切受贿和肮脏之财。\par

\Person{塔拉修斯}，至圣宗主教，说：这整个神圣集会为你们向天主教会所做的这宣认，向天主献上光荣和感谢。\par

神圣会议说：愿光荣归于使分裂者合一的天主。\par

[\Emph{米拉主教\Person{狄奥多尔}接着宣读了同一宣认，并被接纳。下一个请求被接纳的主教宣读如下：} (col. 60)]\par

\Person{德奥多西}，谦卑的基督徒，致神圣大公会议：我宣认并\OriginalTerm{同意}{\GreekText{συντίθεμαι}}并接受并致敬并敬礼，首先，我们的主耶稣基督、我们真实天主的无玷圣像，以及那无玷地生育祂的圣天主之母的神圣圣像，并每天每夜作为一个罪人呼求她的助佑、保护及转求，因为她对基督我们的天主有信赖，因为祂由她所生。同样，我也接受并敬礼神圣而最可赞颂的宗徒、先知、殉道者以及教父和旷野修士的圣像。当然，并非作为神（天主管制！）我以全心求他们所有人为我向天主祈祷，使祂因他们的转求赐我在审判之日得蒙祂的仁慈，因为借此我更能彰显我从起初对他们所有的情感和爱情。同样，我也敬礼并尊敬并致敬圣人的遗骸，作为为基督奋战者，他们从祂领受了医治疾病、痊愈伤痛、驱逐魔鬼的恩宠，正如基督教会从圣宗徒和教父传承至今。\par

此外，我甚愿在信友的教堂中有圣像，尤其我们的主耶稣基督和圣天主之母的圣像，用各种材料，即金、银和各色颜料，以便祂的降生成人昭示于众人。同样，可描绘圣人和先知及殉道者的生平，以便他们的奋斗和痛苦简要地显示出来，以激励和教导百姓，尤其未受过教育者。\par

因为若民众拿着灯和香出去迎接送到城市或乡间的皇帝的桂冠像和画像，他们所尊敬的肯定不是涂蜡的木板，而是皇帝本人。更何况在我们天主基督的教会中，应当描绘天主我们的救主及其无玷之母以及所有圣善而蒙福的教父和苦修者的圣像呢？正如圣巴西略所说：\Quote{作家和画家都描绘战争的伟业；前者用言语，后者用画笔；两者都激励许多人勇于作战。}同样这位作者又说：\Quote{你曾费了多少力气，才能找到一位圣者愿意作你向主恳切转求的中保？}金口若望说：\Quote{圣人逝世后他们的爱德并不减少，也不会随离世而终止，反而死后比在世时更加强有力。}以及其他数不胜数的话。因此我恳求你们，诸位圣人！我向你们呼号。我得罪了天，也得罪了你们。请你们接纳我，如同天主接纳了那个奢华的人、那位妓女和那个盗贼。请寻找我，如同基督寻找那只迷失的羊，他把羊扛在肩上；这样，因你们的转求，在天主和祂的天使面前，将因我的悔改和得救而欢乐，至圣的主们！凡不敬礼圣善可敬的圣像者，绝罚！凡亵渎尊贵可敬圣像者，绝罚！凡胆敢攻击和亵渎可敬圣像并称它们为偶像者，绝罚！对基督信仰的诽谤者，即圣像破坏者，绝罚！凡不殷勤教导所有爱基督的人民敬礼和问候各世代中悦乐天主的众圣人的可敬、神圣而尊贵的圣像者，绝罚！凡心怀疑虑而不全心承认他们敬礼神圣圣像者，绝罚！\par

斯图狄奥斯隐修院最可敬的院长萨巴斯说：按照宗徒的训令和普世公会议，他应被接纳回来。\par

最神圣的宗主教塔拉修斯说：从前是正统信仰的诽谤者，如今成了真理的辩护人。\par

\EditorialNote{本次会议接近尾声时（col. 77）}\par

最可敬的东方大祭司的代表若望主教说：这个异端是所有异端中最坏的。祸哉，圣像破坏者！它是最坏的异端，因为它颠覆了我们救主的降生成人（\OriginalTerm{救恩工程}{\GreekText{οἰκονομίαν}}）。\par

% structure: p0050 source=h3 target=subsection reason=relative-heading-rank; base=h2

\subsection{第二次会议}

\EditorialNote{教宗书信由教宗代表呈上。首先宣读的是致君士坦丁和依勒纳的书信，但并非全文，如果我们能信任图书馆员阿纳斯塔修斯的话，他给出了他所称的拉丁原文。以下是这封信的译本以及希腊文本，还有完全删去的拉丁段落的译文（据称是经罗马教宗代表同意的）。}\par

教宗哈德良书信的一部分。\par

\EditorialNote{教宗所写原文。}\par

\EditorialNote{Migne, \WorkTitle{拉丁教父集}, 第96卷, 1217栏}\par

如果你们持守你们所开始的这个正统信仰，并且通过你们，神圣可敬的圣像在那些地区重新树立起来，如同由虔诚记忆的君主皇帝君士坦丁和真福海伦所行的那样——他们传播了正统信仰，并高举了你们属灵的母亲——圣公教会和宗徒的罗马教会，并与其他正统皇帝一起敬拜她作为所有教会的元首——那么，你们那蒙天主护佑的仁慈将获得另一个君士坦丁和另一个海伦的名号，圣公教会和宗徒的教会从起初藉他们获得了力量，如同他们一样，你们的皇帝声名因胜利而远播，以致在全世界光辉灿烂、深植人心。但更重要的是，如果你们遵循正统信仰的传统，接受真福伯多禄宗徒之长的教会的判断，并且如同古时你们的先辈圣皇帝们所做的那样，你们也以尊荣敬拜他的代理人，并以你们神圣的陛下优先遵循我们圣罗马教会的正统信仰。愿宗徒之长本人——我们的主天主曾赐予他权能在天上去束缚和赦免罪恶——常作你们的保护者，将一切野蛮民族践踏在你们脚下，并使你们在各地得胜。因为神圣的权威揭示了其尊荣的标记，以及所有世上的信友应如何向他的至高宗座表示极大的敬礼。因为主将那执掌天国钥匙的人立为万有之首，并藉此特权使他受尊荣，天国的钥匙托付给了他。因此，这位蒙受如此崇高荣耀的人，被堪当承认那信仰，基督的教会就建立在这信仰之上。真福的宣认随之带来真福的赏报，藉此宣认的宣讲，神圣的普世教会得到了光照，其他天主的教会也从其中获得了信仰的确证。因为真福伯多禄宗徒之长，首先坐在宗徒之座上，将其宗徒职和牧灵关怀的首席权留给了他的继承者们，他们将永远坐在他最神圣的宝座上。而他从主天主我们的救主所领受的权威权能，他也依天主命令赐予并交付给了他的继承者教宗们，等等。\par

\EditorialNote{在公会议上用希腊文宣读。}\par

\EditorialNote{Migne, \WorkTitle{拉丁教父集}, 第96卷, 1218栏}\par

如果古代的正统信仰因你们的助力而在那些地区得以完善和恢复，可敬的圣像得以放回原处，那么你们将有分于如今安息于天主的先帝君士坦丁大帝和皇后赫肋纳，她们曾显明并确认了正统信仰，更加高举了你们的圣教会、公教罗马精神教会；并有分于他们之后统治的正统皇帝们；如此，你们那极其虔诚、蒙天保护的声名本身也将被展示为另一位君士坦丁和另一位赫肋纳，受普世赞扬和称颂，因你们，神圣的公教使徒教会得恢复。尤其如果你们跟随圣伯多禄和保禄——首席宗徒——教会的正统信仰传统，接纳他们的代表，正如你们之前的古代皇帝们曾尊敬并全心热爱他们的代表；如果你们的圣座尊敬至圣罗马教会——首席宗徒的教会——天主的圣言亲自赐予他们权能，在天上和地上释放和束缚罪过。因为他们将为自己的权能展开盾牌，所有蛮族民族都将被置于你们脚下；无论你们到哪里，他们必使你们得胜。因为神圣的首席宗徒们自己建立了公教正统信仰，并且立下文字法律：凡以后继承他们宗座的人，都当持守他们的信仰，并在其中坚持到底。\par

\EditorialNote{\Emph{从未向大公会议宣读的部分。}}\par

\EditorialNote{摘自L. and C.，\Emph{Concilia}，第七卷，第117栏。}\par

我们非常惊讶，在你们颁发给皇城宗主教塔拉修斯的谕令中，我们发现他被称为\OriginalTerm{普世}{Universal}；但我们不知道这是由于无知、分裂，还是恶人的异端。但从今以后，我们劝诫你们最仁慈的皇帝陛下：绝不要在你们的文书中称他为\OriginalTerm{普世}{Universal}，因为这似乎与神圣教规的制度及圣父传统的法令相悖。因为若非我们神圣公教使徒教会的权威，他绝不可能占据第二位，这是众人皆知的。因为如果他被称作\OriginalTerm{普世}{Universal}，高过享有优先地位的神圣罗马教会——她是所有天主教会之首——那么他显然就是反对神圣大公会议的人和异端者。因为，如果他是\OriginalTerm{普世}{Universal}，他就被认为甚至对我们的宗座教会也拥有首席权，这在所有信友看来是荒谬的：因为在全世界，首席地位和权能是由世界的救主亲自赐予有福的宗徒伯多禄的；并通过同一位宗徒——我们卑微地占据他的位置——神圣的公教使徒罗马教会现在和永远拥有首席地位和权能的权威，以致如果有人——我们不信会如此——称他为\OriginalTerm{普世}{Universal}，或同意他被这样称呼，就让他知道他已疏远正统信仰，反对我们神圣的公教使徒教会。\par

\EditorialNote{\Emph{宣读结束后}（第120栏）}\par

至圣宗主教塔拉修斯说：你们自己是否从至圣教宗那里收到这些信件，并将它们带给我们的虔诚皇帝？\par

伯多禄和伯多禄——最敬爱天主的司铎，代表罗马至圣教宗哈德良——说：我们亲自从我们的宗座父那里收到这样的信，并交予虔诚的主上们。\par

最显赫的\OriginalTerm{主计官}{Logothete}若望说：此事也为西西里人所知——敬爱天主的卡塔尼亚主教德奥多禄，以及与他同在的至敬执事埃皮法尼——他代表撒丁岛的总主教。因为这两位奉我们虔诚皇帝之命，与我们的至圣宗主教最可敬的\OriginalTerm{教廷代表}{apocrisarius}一同前往罗马。\par

敬爱天主的卡塔尼亚主教德奥多禄站在中间说：虔诚的皇帝以他尊贵的谕令，派遣至敬爱天主的司铎莱奥（他与我同是您圣德的仆人）携带至圣陛下的宝贵信件；而那位敬重我们的圣德的人——他是我西西里行省的\OriginalTerm{总督}{\GreekText{στρατηγὸς}}——派我带着我们正统皇帝的虔诚谕令前往罗马。\par

当我们到达时，我们宣告了虔诚皇帝的正统信仰。\par

至有福的教宗听后说：既然这事发生在他们的统治时期，天主已高举他们虔诚的统治超越先前所有统治。这篇已宣读的\OriginalTerm{陈情书}{\GreekText{ἀναφορὰν}}连同致您圣德的信函以及在场主持的代表们，他已寄给我们最虔诚的君王们。\par

执事、公证人兼\OriginalTerm{司库}{Cubuclesius}科斯马斯说：还有另一封信由旧罗马至圣教宗寄给我们的至圣普世宗主教塔拉修斯。请按你们神圣会议的决定处置。\par

神圣公会议说：请宣读。\par

\EditorialNote{\Emph{然后宣读了哈德良致君士坦丁堡塔拉修斯的信，信末说：}\Quote{我们亲爱的罗马圣教会首席司铎以及修院院长、司铎兼隐修士伯多禄，已由我们派遣给最安宁虔诚的皇帝们，我们恳求你们为了圣伯多禄——宗徒之长，并为我们的缘故，赐予他们一切仁慈和仁厚的款待，以便为此我们能向你们致以诚挚的感谢。}\Emph{信读完结束时}（第128栏），}\par

最可敬的司铎、旧罗马至圣教宗代表伯多禄和伯多禄说：请至圣的皇城宗主教塔拉修斯说，他是否\OriginalTerm{同意}{\GreekText{στοιχεῖ}}旧罗马至圣教宗的信函？\par

最神圣的宗主教塔拉西乌斯说：充满基督之光、并藉福音生了我们的神圣宗徒保禄，在致罗马人书中称赞他们对我们在基督内真天主所有的真信德的热忱时，这样说：\Quote{你们的信德传遍了全世界。}必须遵循这一见证，反驳它的人缺乏理智。因此，旧罗马的统治者哈德良，既分享这些事，如此被见证，便明确而真实地写信给我们的虔敬皇帝和我们的卑微，极好而优美地确认了天主教会古老传统。我们自己也通过文字、查问、三段论和论证进行了审查，并受教于教父们的教导，如此我们已宣认、正宣认并将宣认；在刚才宣读的书信的意义上，我们将坚定、持守、站稳，接受依照我们圣教父古老传统的肖像描绘；我们以坚定不移的情感敬礼它们，因为它们是因我们天主基督、我们无玷圣母天主之母、圣天使及诸圣之名而制作，最明确地将我们的朝拜和信德归于唯一真天主。\par

神圣公会议说：整个神圣公会议如此教导。\par

彼得与彼得，蒙天主所爱的司铎及宗座代表，说：请神圣公会议说明，是否接受旧罗马最神圣教宗的书信。\par

神圣公会议说：我们追随、接受、承认它们。\par

\EditorialNote{主教们随后逐一表决，均持同一意见。}\par

% structure: p0078 source=h3 target=subsection reason=relative-heading-rank; base=h2

\subsection{第三次会议}

（拉贝与科萨，\WorkTitle{大公会议}，第七卷，第188栏。）\par

君士坦提亚的君士坦丁最神圣主教说：既然我这不配之人发现刚才宣读的、从东方寄给最神圣的总主教暨大公宗主教塔拉西乌斯的书信，与他先前所作的信仰宣认毫无改变，我赞同并同心合意，接受并尊敬地敬礼神圣可敬的圣像。但我将朝拜的敬礼唯独保留给\OriginalTerm{超性}{supersubstantial}而赐予生命的天主圣三。凡不如此想、不如此教导的人，我将其逐出神圣天主教会及宗徒教会，以绝罚击打他们，将他们交付与否认我们的真天主基督降生成人和身体救恩计划之人的命运。\par

% structure: p0081 source=h3 target=subsection reason=relative-heading-rank; base=h2

\subsection{第四次会议}

\EditorialNote{在宣读的众多教父著作中，有一篇取自圣额我略·尼撒的讲道，描述了一幅描绘依撒格祭献的画作，并说明他如何不能\Quote{不流泪}地经过它。}\par

最光荣的君王们说：看我们的教父如何因所描绘的史事而悲伤，甚至流泪。\par

安基拉的最神圣主教巴西略说：教父多次阅读这个故事，但或许未曾流泪；而一旦看见它被描绘出来，他便流泪了。\par

最可敬的隐修士、司铎兼东方高级司祭代表若望说：若图像对如此一位博士有益并引出眼泪，那么在无知和单纯的人身上，它岂不更会带来痛悔和益处？\par

神圣公会议说：我们在多处见过如同教父所说的亚巴郎事迹的绘画。\par

卡塔尼亚的最神圣主教特奥多罗说：若神圣的额我略，在神圣默想中警醒，见到亚巴郎的事迹而感动流泪，那么描绘我们主基督为我们的降生成人的画作，岂不更会感动观看者，使之受益并流泪？\par

最神圣的宗主教塔拉西乌斯说：当我们看见我们被钉之主的图像时，难道不流泪吗？\par

神圣公会议说：我们确实会流泪——因为其中将完美地发现天主为我们的缘故而降生成人的贬抑之深。\par

\EditorialNote{此后不久，宣读了一段圣亚大纳削的著作，描述在贝里图斯发生的奇迹，之后发现以下内容（第224栏）。}\par

最神圣的宗主教塔拉西乌斯说：但或许有人会说：为什么我们所拥有的圣像不行奇迹？对此我们回答：正如宗徒所说，神迹是为不信的人，不是为相信的人。因为走近那圣像的人是不信者。因此天主通过圣像给了他们一个标记，为将他们引向我们的基督徒信仰。但\BibleQuote{邪恶淫乱的世代寻求征兆，但必不给它任何征兆。}\par

\EditorialNote{在宣读了许多其他引文之后，作为第六次公会议的一条教规，宣读了特鲁罗会议的教规（第233栏）。}\par

最神圣的宗主教塔拉西乌斯说：有些人患有无知的病症，对这些教规（即特鲁罗会议的教规）感到惊讶，说：你们真的认为它们是在第六次公会议上采纳的吗？现在让所有这样的人知道：神圣伟大的第六次公会议是在君士坦丁堡召开的，针对那些说基督只有一种能量和意志的人。这些会议谴责了异端者，阐明了正统信仰，然后在君士坦丁在位的第十四年各自回家。但四五年后，同一批教父在君士坦丁之子尤斯蒂尼安下再次聚集，制定了上述教规。对此勿存疑虑。因为在君士坦丁下签名的人，与在尤斯蒂尼安下签署这份文件的人是同一批人，这可以从他们手迹不变的一致性清楚表明。因为凡出现在大公会议上的人，也应当制定教会教规。他们说，我们应当被可敬的圣像引导（如同被手引导）去纪念基督的降生成人及其救恩之死，若藉由它们我们被引向认识我们天主基督的降生成人，那么我们对那些推倒可敬圣像的人会持有什么意见呢？\par

\EditorialNote{在会议结束时，宣布了一系列绝罚后，宣读了以下内容，所有主教都在其上签名（第317栏）。}\par

完成我们天主和救主耶稣基督的神圣诫命，我们神圣的教父们并没有将祂所赐予他们的神圣知识之光藏在斗底下，而是将它放在最有益教导的灯台上，好让它照亮屋里所有的人——也就是说，照亮那些生于天主教会中的人；免得凡虔诚承认主的人中有一位可能将脚撞在异端邪恶教义的石头上。因为他们驱逐了异端者的每一错误，并砍掉了腐烂的肢体（如果它已无可救药地病入膏肓）。他们用簸箕扬净了禾场。好麦子，即那滋养并强壮人心的圣言，他们收进了天主教会的粮仓；而将异端邪说的糠秕扔到外面，用不灭的火焚烧。因此，这第二次聚集在这光荣的\Place{尼西亚}大都会的神圣大公会议，因天主的旨意，并奉我们虔诚而最忠信皇帝的命令——\Person{伊琳娜}，一位新的\Person{海伦}，以及她受天主保护的后代，一位新的\Person{君士坦丁}——通过阅读我们经核准且有福的教父们的教导，已经光荣了天主自己，从祂那里赐予了他们智慧以教诲我们，并为完善天主教会和宗徒教会；而对那些不信从他们（的教导）、却企图通过他们的新奇来遮蔽真理的人，他们诵唱了诗篇的话：\Quote{你的仇敌在你的圣所中行了大恶；并自夸说，不再有师傅了，他们不知道我们诡诈地对待了真理之言。}但我们在一切事上持守我们同样神载教父们的教义和诫命，以同一口吻和同一心灵宣告，既不增加也不减去他们传授给我们的任何东西。我们靠着这些得到坚固，靠着这些得到确认。我们如此宣认，如此教导，正如神圣而大公的六次大公会议所规定和批准的。我们信唯一的天主，全能的圣父，一切有形无形之物的创造者；并信唯一的主耶稣基督，祂的独生子圣言，万物藉祂而受造；并信圣神，主及赋予生命者，与同一圣父和祂无始的圣子\OriginalTerm{同体}{consubstantial}、\OriginalTerm{同永恒}{coeternal}。那未受造的、不可分的、不可理解的、不限定的天主圣三；祂，整个且独一，应受朝拜和敬仰的钦崇；一个神性，一个主性，一个主权，一个王国和王朝，无分地分施于各位置，且分别适合于本质。因为我们宣认同一神圣且同体的天主圣三中的一位，我们的主耶稣基督真天主，在末后的日子为我们的救恩降生成人并成为人，通过祂救恩的降生、苦难、复活和升天拯救了我们的种族；并将我们从偶像的错误中解救出来；正如先知所说：\Quote{不是使者，不是天使，而是主自己拯救了我们。}我们也跟随祂，采纳祂的声音，大声呼喊：\Quote{没有大公会议，没有君王权能，没有天主所憎恨的协议将教会从偶像的错误中解救出来}，正如那犹太化非法会议疯狂妄想的那样，它曾对可敬的圣像狂怒；但荣耀的主自己，降生成人的天主，拯救了我们，将我们从偶像崇拜的欺骗中夺回。因此愿光荣归于祂，愿感谢归于祂，愿感恩祭归于祂，愿赞美归于祂，愿尊威归于祂。因为祂的救赎和祂的救恩本身能完美地拯救，而非其他来自尘土的人。因为祂亲自为我们——我们这些通过祂降生成人的救世奥迹而到了末世的人——实现了祂先知们预先所说的话，因为祂居住在我们中间，在我们中间往来，并从地上除去了偶像的名字，正如所记载的。但我们尊崇主和祂宗徒们的声音，藉此我们被教导：首先要尊敬那确切而真正是天主之母、且高乎一切天上万军者；也尊敬神圣的天使军旅；以及蒙福且备受赞颂的宗徒们、荣耀的先知们、为基督而战的胜利的殉道者们、神圣而神载的圣师们、以及所有圣人；并寻求他们的转求，因为他们能够使我们在全能的天主面前蒙恩，只要我们遵守祂的诫命，并努力行善。此外，我们尊崇可敬而赋予生命的十字架圣像，以及圣人们的圣髑；我们接受神圣而可敬的圣像：我们尊崇、拥抱它们，按照神圣天主教会——即我们神圣教父们——的古老传统，他们也接受了这些并设立在所有最神圣的天主教堂中以及祂统御的每一处。这些可敬而可尊的圣像，如前所述，我们尊敬、尊崇并虔诚地敬礼：即我们伟大的天主和救主耶稣基督降生成人的圣像，以及我们无玷的圣母、全圣的天主之母——祂乐意从她取得肉体，以拯救并将我们从一切不敬偶像崇拜中解救出来——的圣像；也尊敬神圣而无形的天使们的圣像，他们曾以人的形式显现给义人。同样地，神圣而备受赞颂的宗徒们的形像和肖像，以及宣讲天主的先知们、奋斗的殉道者们和圣人们的形像。如此，透过他们的形像，我们得以在记忆和回想中被引向\OriginalTerm{原型}{prototype}，并分享其中某一位的圣德。\par

因此，我们学会了如此思想，并藉由我们圣教父的教导得以坚固，也藉由他们那神圣传授的训诲得以坚强。感谢天主，因祂那不可言喻的恩赐，祂在末时没有抛弃我们，恶人的杖也没有落在义人的份上，免得义人伸手，即他们实际的行为，趋向邪恶。但祂善待那些心地善良正直的人，正如圣咏作者\Person{达味}所歌唱的；我们也与他同唱圣咏的其余部分：\BibleQuote{至于那些转回自己邪恶的人，主必将他们与作恶的人一同带走；愿平安归于天主的以色列。}\par

\EditorialNote{随后是署名，并结束了本次会议的记录（col. 321-346）。}\par

% structure: p0098 source=h3 target=subsection reason=relative-heading-rank; base=h2

\subsection{第六次会议}

（Labbe 与 Cossart，\WorkTitle{大公会议}，第七卷，col. 389。）\par

最著名的书记长\Person{莱奥}说：神圣而可称颂的公会议知道，在上一场会议中，我们如何审查了那些背弃天主的异端者的各种言论，他们曾指控神圣无瑕的基督徒教会设立神圣的图像。但今天，我们手中拿到了这些诽谤基督徒之人的书面亵渎，即那个伪会议荒谬、轻易可驳且自证其罪的\OriginalTerm{定义}{\GreekText{ὅρον}}，它在各方面都与那些为天主所憎恨的异端者的邪恶意见相符。但我们不仅持有此定义，还持有由圣神所督导的精巧且最有力的驳斥。因为应当以明智的反驳来战胜这一定义，并以强有力的驳斥将其粉碎。我们也提交了这些，以便了解诸位对此的意见。\par

神圣的公会议说：请宣读。\par

\Person{若望}，执事兼总务长\EditorialNote{最神圣的君士坦丁堡大教堂的，仅见于拉丁文版本}宣读。\par

\EditorialNote{随后，\Person{若望}（执事）宣读了正统的驳斥，而\Person{额我略}（\Place{新凯撒勒雅}的主教）宣读了伪公会议的定义，一人宣读异端陈述，另一人宣读正统回答。}\par

% structure: p0104 source=h2 target=section reason=relative-heading-rank; base=h2

\section{定义}

神圣而大公的公会议，因天主的恩宠以及蒙天主所爱且正统的皇帝\Person{君士坦丁}与\Person{莱奥}的最虔诚之命令，如今在帝国都城，在神圣无玷的天主之母、童贞玛利亚的圣殿（名为\Place{布拉赫内宫}）中聚会，所颁布如下。\par

撒殚误导了人，以致他们崇拜受造物而非造物主。梅瑟法律和先知们曾协力消除这一败坏；但为了彻底拯救人类，天主派遣了自己的圣子，祂使我们脱离谬误和偶像崇拜，教导我们以心神和真理崇拜天主。作为祂救恩学说的使者，祂给我们留下了宗徒和门徒，他们以祂荣耀的教义装饰了教会——祂的新娘。教会的这一装饰，圣教父和六次大公会议都予以保持无损。但那先前提到的邪恶的工匠不能忍受看到这一装饰，便逐渐以基督教的外表带回了偶像崇拜。正如基督曾以圣神的能力武装祂的宗徒对抗古代的偶像崇拜，并派遣他们到普世，同样，祂也兴起祂的仆人——我们忠信的皇帝——对抗新的偶像崇拜，并以同样的圣神智慧赐予他们。受圣神的推动，他们不能再目睹教会因魔鬼的欺骗而荒废，于是召集了蒙天主所爱的主教们的神圣聚会，以便在会议中依照圣经审查那些图像的欺骗性着色（\OriginalTerm{图像}{\GreekText{ὁμοιωμάτων}}），这些图像将人的精神从对天主崇高的\OriginalTerm{朝拜}{\GreekText{λατρείας}}引向对受造物低劣而物质的\OriginalTerm{朝拜}{\GreekText{λατρείαν}}，并期望他们能在神的引导下对此发表意见。\par

因此，我们的神圣公会议聚集了，我们这三百三十八位成员遵循先前的会议法令，接受并欣然宣告所传授的信理，尤其是六次神圣大公会议的信理。首先是神圣而大公的尼西亚大公会议，等等。\par

在圣神的引导下仔细审查了他们的法令之后，我们发现那描绘活物的非法艺术亵渎了我们救恩的基本信理——即基督的降生成人——并且与六次神圣公会议相矛盾。这些会议谴责了\Person{聂斯多略}，因为他将天主的独一圣子和圣言分为两个儿子；另一方面，他们也谴责了\Person{亚略}、\Person{狄约斯哥}、\Person{欧迪奇}和\Person{塞维鲁}，因为他们坚持基督二性的混合。\par

因此，我们认为应当在我们目前的定义中，极其准确地揭示那些制作并敬拜此类像者的谬误，因为所有圣教父以及六次大公会议一致教导：任何人不得想象在一位位格或一位之内两性结合时存在任何分离或混合以反对那不可探究、不可言说、不可理解的结合。那么，那因贪财的罪恶之爱而描绘那不该描绘之物的画师的愚蠢有何益处？也就是说，他用自己污秽的手试图塑造那应当只在心中相信、用口承认的事物。他制造了一幅像，并称之为基督。\Quote{基督}之名意指天主和人。因此，那是天主和人的像，从而他在自己愚昧的心思中，在受造肉身的再现中，描绘了那不可描绘的天主性，从而混合了不可混合之物。这样，他犯有双重亵渎：一是制作天主性的像，二是混合天主性与人性。那些敬拜这像的人也陷入同样的亵渎，两者都承受同样的灾祸，因为他们与亚略、狄奥斯哥鲁、欧迪奇一同犯错，并与无首派的异端相同。然而，当他们被责备竟敢描绘那不可描绘的基督天主性时，他们以这样借口逃避：我们只描绘我们所看见、所摸过的基督的肉性。但那是涅斯多留派的错误。因为应当考虑：那肉性也是天主圣言的肉性，毫无分离，被天主性完全摄取，并完全成为神圣。如今它怎能被分开并单独描绘呢？基督的人性灵魂也是如此，它在圣子的天主性与肉体的迟钝中间作中保。就如人肉身同时是天主圣言的肉身，同样人灵魂也是天主圣言的灵魂，两者同时存在，灵魂与身体都被神化，天主性在祂自愿受难中，甚至在灵魂与身体分离时，也保持不可分。因为基督的灵魂在哪里，祂的天主性也在那里；基督的身体在哪里，祂的天主性也在那里。那么，如果在祂受难时天主性从未与这些分离，这些愚人怎敢将肉性与天主性分开，并单独把肉性作为纯人的像来描绘？他们陷入不敬虔的深渊，因为将肉性与天主性分开，赋予肉性自有的存在、自有的位格，并将其描绘，从而在天主圣三中引入第四位。此外，他们把那因被天主性摄取而成为神圣的，当作未被神化的来描绘。因此，凡制作基督像的，要么描绘那不可描绘的天主性，并将其与人性混合（如一性论者），要么描绘基督的身体为未神化的、分离的、独立的位格，如涅斯多留派。\par

然而，基督人性唯一可容许的形象，是圣体圣事中的饼和酒。祂选择了这一形式，而非其他形式，这一预像，而非其他预像，来代表祂的降生成人。祂命令献上饼，而不是人形的肖像，以免偶像崇拜出现。并且，就如基督的身体被神化，同样，这基督身体的形象——饼，也因圣神的降临而被神化；它藉司铎的中介成为基督的神圣身体，司铎将祭品与凡俗之物分离，并祝圣它。\par

给图像命名的恶习并非来自基督、宗徒和圣教父；他们也没有留下任何祷文，藉以祝圣图像，或使其超越普通物质。\par

然而，若有人说，我们关于基督像的论述可能是正确的，因为两性神秘的结合，但禁止那全然无玷、永荣的天主之母、先知、宗徒和殉道者的像——他们只是人，并非由两性组成——则是不对的；我们可以首先回答：如果那些像被弃绝，这些也就不必要了。但我们也要考虑针对这些特别可能提出的反驳。基督教拒绝了整个异教，因此不仅拒绝异教祭献，也拒绝异教偶像崇拜。圣人们虽已死去，却与天主永远同活。若有人想通过一种死的艺术——异教所发明的手段——将他们召回生命，他就犯有亵渎之罪。谁敢用异教艺术去描绘那超越诸天及诸圣人的天主之母？基督徒怀着复活的希望，不允许效仿崇拜魔鬼者的习俗，用普通死物去侮辱那闪耀于如此荣光中的圣人们。\par

此外，我们可以用圣经和教父来证明我们的观点。圣经中说：\BibleQuote{天主是神：朝拜祂的人应当以心神以真理去朝拜祂。} 以及：\BibleQuote{你们不可为自己雕刻偶像，也不可制造任何天上、地上或地下水中之物的肖像。} 因此天主在火中于山上向以色列人说话，却未向他们显示任何形象。再者：\BibleQuote{他们将不可朽坏的天主的光荣变成了朽坏的人的肖像……反而事奉受造物超过事奉造物主。} \EditorialNote{还引用了其他几处更不贴切的段落。}\par

圣教父们也教导同样的事。\EditorialNote{会议引用了一篇出自圣埃皮法尼乌斯的伪作，以及一篇插入安西拉的德奥多图斯（圣济利禄的朋友）著作中的文字；引用了纳西盎的额我略、圣金口若望、圣巴西略、安菲洛基奥的阿塔纳西乌斯、以及欧瑟伯·庞菲里的言论——这些言论毫无新意——来自他的\WorkTitle{致君士坦莎皇后的信}，她曾向他索要一幅基督画像。}\par

受圣经和教父的支持，我们以天主圣三之名一致声明：任何由画师的恶艺用任何材料和颜色制造的形象，都应在基督教会中被拒绝、移除并诅咒。\par

今后凡胆敢制作、敬礼此类物件，或在圣堂、私宅中摆设，或秘密持有者，若是主教、长老或执事，应被罢免；若是隐修士或平信徒，应受绝罚，并应受世俗法律审判，作为天主的敌人和教父所传教义的仇敌。同时，我们规定：任何圣堂负责人不得以销毁有关图像的谬误为借口，动手改制圣器，因其上饰有图像。对圣堂祭衣、布品及一切奉献于神圣礼仪之物亦作同样规定。然而，若有圣堂负责人希望改制此类圣器及祭衣，他必须征得神圣普世宗主教同意并遵奉我们虔诚皇帝的谕令方可进行。同样，任何王侯或世俗官员不得掠夺圣堂，如以往有些人以销毁图像为借口所为。我们制定这一切，相信我们是如同宗徒一样说话，因为我们也相信我们拥有基督的圣神；正如我们信德相同的前辈们按主教会议制定的内容发言，我们也相信而如此说话，并制定一项定义，表明我们认为合宜的，是遵循并符合我们教父们的定义。\par

(1) 任何人若不按宗徒及教父的传统，宣认在圣父、圣子及圣神内，一个天主性、本性及实体，意志及行动，德能与统御，国度与权能，存在于三个自立体，即他们最荣耀的位格中，则受绝罚。\par

(2) 任何人若不宣认天主圣三之一位成了血肉，则受绝罚。\par

(3) 任何人若不宣认童贞玛利亚确实是天主之母等，则受绝罚。\par

(4) 任何人若不宣认一个基督，既是天主又是人，则受绝罚。\par

(5) 任何人若不宣认主的肉体是赋予生命的，因为它是天主圣言的肉体，则受绝罚。\par

(6) 任何人若不宣认基督内有两个本性，则受绝罚。\par

(7) 任何人若不宣认基督带着肉体及灵魂与天主圣父同坐，并将如此来审判，且他将永远为天主，毫无粗鄙性等，则受绝罚。\par

(8) 任何人若胆敢用物质颜料描绘圣言降生成人后的神圣\OriginalTerm{形象}{\GreekText{χαρακτήρ}}，则受绝罚！\par

(9) 任何人若胆敢因降生成人而用物质颜料以人形描绘不可描绘的圣言之本体或位格\OriginalTerm{本体}{ousia}或\OriginalTerm{位格}{hypostasis}，而不更宣认即使在降生成人之后他（即圣言）仍是不可描绘的，则受绝罚！\par

(10) 任何人若胆敢在图画中描绘两个本性的位格性结合，并称之为基督，从而错误地表现两个本性的结合等，则受绝罚！\par

(11) 任何人若将与圣言之位格结合的肉体与圣言分离，并试图在图画中单独描绘它等，则受绝罚！\par

(12) 任何人若将唯一基督分成两个位格，并试图单独描绘由童贞玛利亚所生者，从而只承认诸本性的相对\OriginalTerm{相对的}{\GreekText{σχετική}}结合等，则受绝罚。\par

(13) 任何人若在图画中描绘因与圣言结合而被神化的肉体，从而将其与天主性分离等，则受绝罚。\par

(14) 任何人若试图用物质颜料将天主圣言描绘为纯粹的人，虽然他具有天主的形状，却以自身位格取了奴仆的形状，从而试图将他与不可分离的天主性分开，以致在天主圣三中引入了四位一体等，则受绝罚。\par

(15) 任何人若不宣认神圣、卒世童贞玛利亚，真实且恰当的天主之母，高于一切受造物，无论有形无形，且不以真诚的信德寻求她的转求，如同寻求一位有自信能接近我们天主者，因为她生了祂等，则受绝罚。\par

(16) 任何人若试图用毫无价值的物质颜料在无生命的图画中描绘圣人们的形象（因为此思想是虚妄的，由魔鬼引入），而不更在自身中描绘他们的德行作为活生生的肖像等，则受绝罚。\par

(17) 任何人若否认呼求圣人的益处等，则受绝罚。\par

(18) 任何人若否认死者的复活、审判、各人应得的报应、无尽的痛苦和无尽的福乐等，则受绝罚。\par

(19) 任何人若不接受此我们神圣大公第七次大公会议，则从圣父、圣子及圣神，并从七次神圣大公会议受绝罚！\par

\Emph{[随后是针对制定或传授任何其他信仰的禁令，以及不服从的惩罚。此后是欢呼赞美。]}\par

神圣皇帝君士坦丁与良说道：愿神圣大公会议表明，是否在全体至圣主教同意下，刚才宣读的定义已经颁布。\par

神圣会议喊道：\Quote{我们全都如此相信，我们全都同心同意。我们全都异口同声并自愿签署。这是宗徒们的信德。愿皇帝万岁！他们是正统信仰之光！愿正统皇帝万岁！愿天主护佑你们的帝国！你们如今更坚定地宣告了基督两个本性的不可分割！你们驱逐了所有偶像崇拜！你们摧毁了（君士坦丁堡的）哲曼努斯、乔治及曼苏尔\OriginalTerm{曼苏尔}{\GreekText{μανσουρ}}（若望·达玛森）的异端。绝罚哲曼努斯，这心怀二意者、崇拜木头者！绝罚乔治，他的同伙，教父教义的篡改者！绝罚曼苏尔，他有恶名和撒拉森人的见解！绝罚背叛基督的人、帝国之敌、不虔敬之师、圣经曲解者曼苏尔！天主圣三已废黜了这三者！}\par

% structure: p0139 source=h2 target=section reason=relative-heading-rank; base=h2

\section{法令}

(\Emph{见拉贝与科萨}，\WorkTitle{会议录}。\Emph{第七卷，第552栏}。)\par

神圣、伟大、大公会议，因天主的恩宠及虔诚爱基督的皇帝君士坦丁与其母依琳娜的旨意，在比提尼亚辉煌的都会尼西亚第二次召开，于被称为索菲亚的天主圣堂中，遵循了天主教会，定义如下：\par

我主基督曾将认识祂自己的光明恩赐给我们，并从崇拜偶像的疯狂黑暗中将我们救赎，祂娶了至圣天主教会，毫无玷污或缺陷，并应许要如此保存她：祂向圣门徒们说了这话：\BibleQuote{看！我天天与你们同在，直到世界的终结。}\BibleRef{玛 28:20}这应许不只为他们，也为我们这些因他们的话而相信祂名的人。但有些人，不考虑这份恩赐，因狡猾仇敌的诱惑而变得反复无常，从正信上堕落了；因为他们背离了天主教的传统，在真理上误入歧途，如同谚语所说：\Quote{农夫在自己的耕作中走迷了路，手中收集了虚无。}因为某些司铎——仅是名义上的司铎，而非实际上的——竟敢反对圣所中天主所悦纳的装饰；关于他们，天主借先知大声说：\BibleQuote{许多牧人毁坏了我的葡萄园，玷污了我的产业。}\BibleRef{耶 12:10}\par

确实，他们追随亵圣之人，被肉性知觉引入歧途，诽谤我们天主基督的教会——就是祂所娶的——并且不能区分圣与俗，将我们主及祂的圣人们的肖像与魔鬼偶像的雕像等同视之。我们的主天主看见这些事（不愿看到祂的子民被如此瘟疫所腐化），乐意将我们——祂的主要司铎们——从各方召集，怀着神圣的热忱，并因我们的君主君士坦丁与伊琳的旨意而来到这里，为使天主教的传统能借我们的共同法令得以稳固。因此，我们尽心竭力，经过透彻的审查和分析，遵循真理的方向，不减少，不增添，但完好保存一切属于天主教会的事，并追随六大公会议，尤其在此辉煌的尼西亚大都会召开的会议，以及后来在受天主保护的皇城所召集的会议。\par

我们信……来世的生命。阿们。\par

我们憎恶并弃绝亚略和所有与他荒谬意见相同的人；也弃绝马其顿纽及其追随者——他们被恰当地称为\Quote{\OriginalTerm{圣神敌对者}{Pneumatomachi}}。我们宣认我们的圣母、圣玛利亚确实且真实地是天主之母，因为她在肉身方面是天主圣三一位——即我们天主基督——的母亲，正如厄弗所大公会议所定义，当时它将不敬的聂斯多略及其同伙逐出教会，因为他教导〔基督内〕有两个位格。与本届会议的教父们一起，我们宣认：由无玷天主之母、卒世童贞玛利亚降生成人的那一位，有两个性体，承认祂是完美天主和完美的人，正如加采东大公会议所颁布，将欧提克斯和狄奥斯库若作为亵渎者逐出神圣的庭院\OriginalTerm{庭院}{\GreekText{αὐλῆς}}；并将塞维鲁、伯多禄及其他以各种方式亵渎的人列入同一行列。此外，与这些一起，我们弃绝奥力振、厄瓦格利乌与狄迪莫的虚构说法，依照在君士坦丁堡举行的第五次大公会议的决定。我们肯定在基督内有两个意志和两个行动，符合每个性的实际，正如在君士坦丁堡举行的第六次大公会议所教导，将色尔爵、何诺留、居鲁士、丕洛、玛加略及其同意他们的人，以及所有不愿恭敬的人逐出。\par

简而言之，我们完整保存传授给我们的所有教会传统，无论是以文字或口传的方式，其中之一便是绘制图像，这符合福音宣讲的历史，是一种在许多方面有益的传统，尤其在此：即如此，天主圣言的降生成人被显示为真实而非仅仅是幻想的，因为这些（图像和福音）彼此指示，并且无疑也具有相互的意指。\par

因此，我们遵循王道和圣教父们受天主启发的权威以及天主教的传统（因为，众所周知，圣神住在教会内），准确无误地规定：正如珍贵而赋予生命的十字架的形状，同样地，可敬而神圣的肖像，无论以绘画、镶嵌或其他合适材料制成，都应放在天主的圣堂中，在圣器、祭衣、挂幔上，以及在画像中，无论是在家中或路边，即：我们主天主救主耶稣基督的肖像、我们无玷的圣母天主之母的肖像、可敬天使的肖像、所有圣人和所有虔诚者的肖像。因为它们在艺术表现中被看到的次数越多，人们就越容易被提升至对其原型的记忆和渴望；应给予它们合乎礼仪的敬礼和尊崇\OriginalTerm{敬礼与尊崇}{\GreekText{ἀσπασμὸν} \GreekText{καὶ} \GreekText{τιμητικὴν} \GreekText{προσκύνησιν}}，而非那种真正信德的朝拜\OriginalTerm{真正信德的朝拜}{\GreekText{λατρείαν}}——这只能归于神圣本性；但对于这些，如同对珍贵而赋予生命的十字架的形状、以及福音书和其他圣物，可以依照古代虔诚的习俗献上香和灯。因为给予肖像的荣誉传递到肖像所代表者，而且尊敬肖像的人在其中尊敬所代表的主题。因为这样，我们圣教父们的教导，即从地极到地极领受福音的天主教会的传统，得以坚固。这样，我们跟随在基督内说话的保禄，以及整个神圣宗徒团体和圣教父们，坚持我们所领受的传统。因此，我们预言式地咏唱教会的凯旋颂歌：\BibleQuote{熙雍女子，欢呼喜乐吧！耶路撒冷女子，欢呼吧！全心喜乐欢跃吧！上主已从你身上除去仇敌的压迫；你从你敌人的手中被救赎。上主在你中间是君王；你不会再见邪恶，平安与你永远同在。}\BibleRef{匝 9:9-10}\par

因此，凡胆敢作如是思想或教导，或如同邪恶异端者，摒弃教会的传统并发明新奇事物，或拒绝教会所接受的某些事物（\Emph{例如}，福音书、十字架圣像、绘画圣像，或殉道者的神圣遗髑），或邪恶而刻薄地策划任何颠覆天主教会合法传统之事，或将圣器或可敬的隐修院转为俗用，若他们是主教或圣职人员，我们命令他们应被撤职；若是修道者或平信徒，则应断绝共融。\par

[当所有人签署后，欢呼开始（col. 576）。]\par

神圣会议高呼：我们众人都相信，众人都如此心意，众人都同意并已签署。这是宗徒的信仰，这是正教徒的信仰，这是使全世界坚固的信仰。我们信唯一的天主，应受天主圣三的颂扬，我们敬礼可敬的圣像！凡不如此持守者，愿其受绝罚。凡不如此思想者，愿其被逐出教会。因为我们追随天主教会最古老的法律。我们遵守教父的法令。我们绝罚那些在天主教会中添加或删减的人。我们绝罚诋毁基督徒者所引入的新奇之物。我们敬礼可敬的圣像。我们绝罚不这样做的人。绝罚那些胆敢将圣经中关于偶像的话应用于可敬圣像的人。绝罚那些不敬礼神圣可敬圣像的人。绝罚那些称神圣圣像为偶像的人。绝罚那些说基督徒将神圣圣像当作神明的人。绝罚那些说除我们的天主基督之外，另有谁救我们脱离偶像的人。绝罚那些胆敢说天主教会曾接受偶像的人。\par

祝皇帝万岁，等等，等等。\par

% structure: p0152 source=h2 target=section reason=relative-heading-rank; base=h2

\section{法令集}

% structure: p0153 source=h3 target=subsection reason=relative-heading-rank; base=h2

\subsection{第一法令}

\Emph{应全面遵守神圣法令。}\par

领受司铎尊位者的典范，见于典章宪令中所载的证言与训诲，我们欣然接受这些，并借用受天主启示的达味的话，向主天主歌唱：\BibleQuote{我喜爱祢法度的道路，如同喜爱一切财富。} \BibleQuote{祢命令的公义，是祢的法度，直到永远。} \BibleQuote{求祢赐我明悟，我便得以存活。} 既然预言的话命令我们永远持守天主的法度并赖之而生，显然它们必须坚定不移、不可更改。因此天主的先知梅瑟如此说：\BibleQuote{不可加添，也不可删减。} 神圣宗徒以它们为荣，喊道：\BibleQuote{连天使也愿意详细察看这些事。} 并且：\BibleQuote{若有人给你们宣讲福音，与你们所领受的不同，当受绝罚。} 既然这些事如此，我们既有如此确凿的见证，便像得了很多战利品的人一样欢喜，欣然拥抱神圣法令，持守其全部规条，完整不变，无论它们是由圣神的号角、著名的宗徒们，或由六次大公会议，或由为颁布上述大公会议法令而召开的地方会议，或由我们的圣教父所制定。因为这一切都由同一圣神光照，制定了适宜的事。因此，凡他们绝罚的，我们也同样绝罚；凡他们撤职的，我们也撤职；凡他们开除教籍的，我们也开除教籍；凡他们交付惩罚的，我们也处以同样的惩罚。而今，\BibleQuote{你们的生活应没有贪图。} 这是被提升到第三层天、听见不可言说之话的神圣宗徒保禄所喊的。\par

\Emph{古代摘要：我们欣然接受神圣法令，即：圣宗徒的法令、六次大公会议的法令，以及地方会议和我们圣教父的法令，如同一圣神所默感。凡他们绝罚的，我们也绝罚；凡他们撤职的，我们也撤职；凡他们开除的，我们也开除；凡他们处以惩罚的，我们也如此处以惩罚。}\par

% structure: p0157 source=h3 target=subsection reason=relative-heading-rank; base=h2

\subsection{第二法令}

\Emph{凡被祝圣为主教者，必须坚定决心遵守法令，否则不得被祝圣。}\par

当我们诵念圣咏集时，我们向天主许诺：\BibleQuote{我要默想祢的律例，不忘记祢的话语。} 所有基督徒遵守此点乃有益之事，但尤其对于领受司铎尊位者。因此我们规定：凡被提升至主教品位者，必须熟记圣咏集，以便从中劝诫和训导其属下所有圣职人员。都主教应认真查考他是否勤勉研读，且不仅偶尔，而是恒常研读神圣法令、福音、神圣宗徒书信集，以及一切其他神圣圣经；是否按天主的诫命生活，并同样教导其子民。因为我们最高司铎职的特殊宝藏（\OriginalTerm{本质}{\GreekText{οὐσία}}）是由天主传授给我们的神谕，即对神圣圣经的真正知识，正如大德尼削所说。若他的心不坚定、甚至不乐意如此行事和教导，则不得被祝圣。因为天主借先知说：\BibleQuote{你抛弃了知识，我也必抛弃你，使你不再作我的司祭。}\par

\Emph{古代摘要：凡将成为主教者，必须熟记圣咏集：他必须透彻理解所读内容，而非仅肤浅地，而是以勤勉关注，即神圣法令、福音、宗徒书信和全部神圣圣经。若没有这样的知识，则不得被祝圣。}\par

% structure: p0161 source=h3 target=subsection reason=relative-heading-rank; base=h2

\subsection{第三法令}

\Emph{拣选主教不属于君主职权。}\par

凡由世俗统治者任命主教、司铎或执事之选举，均属无效，按教令所言：若有主教利用世俗权力，藉此获得任何教会之管辖权，应被罢免，并与所有与他共融者一同被逐出教会。因为被提升至主教职者，必须由主教选举，正如尼西亚圣教父们所颁布的教令所言：主教最好由省内全体主教祝圣；但若因紧急需要或路途遥远而难以安排，则至少三位主教聚集并投票，缺席者亦须以书信表示同意，方可举行祝圣。如此完成之事，在各省内应由其都主教予以确认。\par

\Emph{古代摘要：凡由世俗官长所作之选举均属无效。}\par

% structure: p0165 source=h3 target=subsection reason=relative-heading-rank; base=h2

\subsection{第四条}

\Emph{主教应戒绝收受一切礼物。}\par

教会之使者、神圣宗徒保禄，为厄弗所的司铎们，也为全体司祭职，订立了一条\OriginalTerm{规则}{\GreekText{κανόνα}}，清楚明言：\BibleQuote{我没有贪图过任何人的金银或衣服。我凡事给你们做了榜样，告诉你们必须这样劳动，扶助病弱者；}因为他认为施予比领受更为有福。因此，我们受他教导而决定，在任何情况下，主教不得为不义之财而捏造虚假的罪过借口，向隶属他的主教、司铎或隐修士索要金银或其他礼物。因为宗徒说：\BibleQuote{不义的人不得承受天主的国，}又说：\BibleQuote{儿女不该为父母积蓄，而是父母该为儿女积蓄。}所以，若有人被发现为索要金银或其他礼物，或因个人情绪，而暂停任何隶属其管辖的圣职人员的职务，甚至将其逐出教会，或关闭任何可敬的圣殿，致使天主的礼仪不能在其中举行，将自己的疯狂发泄于无灵之物，显明自己缺乏理智，则他应受自己为他人所设之同样惩罚，他的烦恼将归到他自己头上，成为违反天主诫命和宗徒训诲的人。因为宗徒的\OriginalTerm{至高首领}{\GreekText{ἡ} \GreekText{κερυφαία} \GreekText{ἀκρότης}}\Person{伯多禄}命令说：\BibleQuote{你们要牧放天主托付给你们的羊群，尽监督之责，不是出于不得已，而是出于自愿；不是为不义之财，而是出于热忱；不是做\OriginalTerm{圣职人员}{\GreekText{τῶν} \GreekText{κλήρων}}的主宰，而是做羊群的模范。当总司牧显现时，你们将获得那不朽的荣耀冠冕。}\par

\Emph{古代摘要：我们决定，任何主教不得向隶属其管辖的主教、圣职人员或隐修士勒索金银或其他物品。若有人因金钱或任何其他事物之权势，或凭个人意愿，阻止其属下任何圣职人员举行神圣礼仪，或关闭可敬的圣殿，以致天主的圣事不能在其中举行，他应受到同态报复。因为宗徒伯多禄说：你们要牧放天主的羊群，不是出于不得已，而是出于自愿，并按照天主；不是为不义之财，而是出于热忱；不是做圣职人员的主宰，而是做羊群的模范。}\par

% structure: p0169 source=h3 target=subsection reason=relative-heading-rank; base=h2

\subsection{第五条}

\Emph{凡因圣职人员在教会内被祝圣时未带礼物而对他们加以侮辱者，应处以罚款。}\par

那些怙恶不悛、执迷于罪者，所犯的是至于死的罪；但那些骄傲自高、反对虔敬与真诚，视钱财比服从天主更宝贵，且漠视祂的典章规诫之人，罪更深重。主天主不与他们同在，除非他们或许因自己的跌倒而谦卑，恢复清醒的心智。他们更应以痛悔之心归向天主，祈求如此重罪的赦免与宽恕，不再以不圣洁的馈赠自夸。因为主亲近那些痛悔的人。因此，对于那班因自己的金钱施舍而被按立在教会中而自夸，并凭恃这邪恶陋习（如此远离天主，与整个司祭职不符），傲慢地张开口，以辱骂之言诽谤那些因生活严谨而被圣神拣选、且未凭任何金钱馈赠而被祝圣的人，我们首先规定他们在自己的品级中居末位；但若他们不悔改，则处以罚金。但如果发现任何人曾这样做（即给予金钱）作为圣职授受的代价，应按宗徒法令处理，该法令说：\Quote{若有主教以金钱换取其尊位（对司铎、执事亦然），应予以撤职，并连同其授职者一同撤职，且应完全被断绝共融，如同西满术士被我伯多禄所做的一样。}与此相同的，是我们圣教父们在加采东所订的第二条法令，它说：\Quote{若有人主教以金钱授圣职，出卖那不可出卖之物，并以金钱祝圣主教、乡村主教、司铎、执事，或任何其他列入圣职人员；或为金钱任命任何人担任经济、辩护人或看守之职；或总而言之，为贪图不义之财而做任何违反法令之事——凡试图做此类事者，一经证实，将面临丧失其品级的危险。而被祝圣者不得从这种交易得来的圣职或擢升中获得任何益处；他应保持为凭金钱所得之职衔与责任的外人。若有人出现在如此可耻而渎神的交易中充当中间人，他若是圣职人员，应被革除品级；若是平信徒或隐修士，应被逐出教会。}\par

\Emph{古代摘要：凡以自己因对教会慷慨施金而获得地位为荣，并藐视那些因德行被选、未经任何馈赠而被任命之人，似乎应在其品级中居末位。若他们坚持其道，应受惩罚。但那些为获得圣职而如此施赠者，应被逐出共融，如同西满被伯多禄逐出一样。}\par

% structure: p0173 source=h3 target=subsection reason=relative-heading-rank; base=h2

\subsection{第六条法令}

\Emph{关于在指定时间举行地方会议}\par

既然有一条法令规定，每年两次在各教省举行主教会议，进行法规调查；但鉴于与会者旅途不便，第六届大公会议的圣教父们规定，每年一次，不论地点或任何托辞，应举行会议，纠正不当之事。我们现在重申这条法令。若发现任何执政者阻挠此事，应被逐出教会。但若任何都主教在无强制、恐惧或其他合理借口的情况下，不关心此事，应受法规处罚。当会议审议法令或有关福音之事时，与会主教们应以全神贯注和严肃思考，守护天主神圣而赐生命的诫命，因为遵守它们有极大的赏报；因为诫命是灯，法律是光，考验和管教是生命之路，主的诫命远照，照亮眼睛。都主教不得向主教们索要他们带来的任何物品，无论是马匹或其他礼物。若他被证实做了此类事，应偿还四倍。\par

\Emph{古代摘要：若不能按先前制定、每年两次举行会议的规定，则至少应如第六届大公会议所定，每年举行一次。若有官员禁止此会，应被逐出；主教若不尽力召集会议，应受惩罚。而当会议举行时，若都主教被发现有从任何主教那里取走物品，应偿还四倍。}\par

% structure: p0177 source=h3 target=subsection reason=relative-heading-rank; base=h2

\subsection{第七条法令}

\Emph{关于未安置圣人圣髑而祝圣的教堂，应补足此缺憾}\par

圣宗徒保禄说：\BibleQuote{有些人的罪是明显的，如同先到审判台前；有些人的罪是随后跟来的。}这些是他们的首要之罪，其他罪则跟随其后。因此，紧随诬蔑基督徒者的异端之后，其他不虔敬之事接踵而来。因为他们从教堂中移除了可敬圣像的临在，同样也抛弃了其他习俗，我们现在必须以成文和不成文的法律予以复兴和维持。因此，我们规定，凡未安置殉道者圣髑而被祝圣的圣殿，应以惯常礼仪安放圣髑。若今后发现有主教在未安置圣髑的情况下祝圣圣殿，应被撤职，作为教会传统的违反者。\par

\Emph{古代摘要：应在那些未安置圣髑而祝圣的教堂中安置圣殉道者的圣髑，并附以惯常祈祷。但任何人若未安置圣髑而祝圣教堂，应被撤职，作为教会传统的违反者。}\par

% structure: p0181 source=h3 target=subsection reason=relative-heading-rank; base=h2

\subsection{第八条法令}

\Emph{关于希伯来人：除非诚心悔改，不应接纳}\par

鉴于某些人固守希伯来人的迷信，竟敢嘲弄我们的天主基督，并假装皈依基督宗教而否认祂，私下并秘密地守安息日及遵守其他犹太习俗，我们决定：这样的人不得参与共融，不得参与祈祷，亦不得进入教会；而应按其宗教公开作为希伯来人，且不得为他们的子女施行圣洗，也不得购买或拥有奴隶。但若其中有任何人，出于真诚之心并怀着信德归化，且全心宣认，摒弃他们的习俗和规条，并为使他人信服而归化，则此人应被接纳并受洗，其子女亦然；且应教导他们谨慎远离希伯来人的规条。但若他们不愿如此行，则决不可被接纳。\par

\Emph{古代摘要：不得接纳希伯来人，除非他们显然以真诚之心归化。}\par

% structure: p0185 source=h3 target=subsection reason=relative-heading-rank; base=h2

\subsection{第九条}

\Emph{不得隐藏任何包含诬蔑基督徒之异端的书籍。}\par

所有针对可敬圣像谬写的一切幼稚虚构和疯狂妄言，必须交给君士坦丁堡主教府，以便与其他异端书籍一同锁藏。若有人被发现隐藏此类书籍，如是主教、司铎或执事，应被罢免；若是隐修士或平信徒，应受\OriginalTerm{绝罚}{\GreekText{ἀνάθεμα}}。\par

\Emph{古代摘要：若有人被发现隐藏针对可敬圣像所写的书，若属圣职人员，应被罢免；若是平信徒或隐修士，应被绝罚。}\par

% structure: p0189 source=h3 target=subsection reason=relative-heading-rank; base=h2

\subsection{第十条}

\Emph{任何圣职人员未经主教知晓，不得离开其教区前往另一教区。}\par

鉴于某些圣职人员误解教会法规，离开自己的教区，奔赴其他教区，尤其是来到这座蒙天主保护的皇城，并寄居于权贵之家，在其祈祷所举行礼仪，故不准许未经其本主教及君士坦丁堡主教许可，将此类人接纳入任何房屋或教堂。若有任何圣职人员未经此许可而如此行，并坚持如此，应被罢免。至于那些在以上各位主教知晓下如此行者，他们不可承担世俗及现世的责任，因为神圣法规禁止此事。若有人被发现担任所谓\OriginalTerm{高级官员}{\GreekText{Μειζότεροι}}之职，则他应放弃此职，否则应从司祭职中被罢免。他更应作儿童及家中其他人的教导者，向他们诵读圣经，因为他领受司祭职正是为此。\par

\Emph{古代摘要：圣职人员离开本堂区后定居于远离其本主教及君士坦丁堡主教之处，不得被接纳入房屋或教堂。若他固执己行，应被罢免。但若他们在知晓我们所说之事的情况下如此行，则不得接受世俗职位；若已接受，应放弃。若拒绝，应被罢免。}\par

% structure: p0193 source=h3 target=subsection reason=relative-heading-rank; base=h2

\subsection{第十一条}

\Emph{应在主教府和隐修院设立\OriginalTerm{经济管理员们}{\GreekText{Οἰκονόμοι}}。}\par

既然我们有义务遵守所有神圣法规，我们务必完整维护那条规定各教会应有\OriginalTerm{经济管理员}{\GreekText{οἰκονόμοι}}的规定。若都主教在其教会任命了一位\OriginalTerm{经济管理员}{\GreekText{οἰκονόμος}}，他做得对；但若他不任命，则允许君士坦丁堡主教凭其\OriginalTerm{自己的}{\GreekText{ἰδίας}}权柄为都主教的教会选定一位经济管理员。若归属他们的主教不愿在其教会任命经济管理员们，都主教亦有类似权柄。隐修院也应遵守同一规则。\par

\Emph{古代摘要：若都主教不选举都主教座堂的经济管理员，宗主教应如此行。若主教不如此行，都主教应如此行；因为这样令加采东集会的教父们满意。隐修院也应遵守同一法律。}\par

% structure: p0197 source=h3 target=subsection reason=relative-heading-rank; base=h2

\subsection{第十二条}

\Emph{主教或\OriginalTerm{院长}{\GreekText{ἡγούμενος}}不得转让教会庄园的任何部分。}\par

若发现主教或\OriginalTerm{院长}{\GreekText{ἡγούμενος}}将主教区或隐修院田产的任何部分转让给世俗权贵，或交给任何其他人，此行为依神圣宗徒的法规无效，该法规说：\Quote{让主教照管教会的一切财物，并如同在天主面前一样管理之。}他不可为自己擅取任何部分，也不可将属于天主之物赐予其亲属。若他们贫穷，应在穷人中被帮助；但不可将其用作盗取教会财产的借口。若有人声称土地只带来损失而无收益，也不可将该地给予邻近的世俗统治者；而应交给圣职人员或农夫。若他们使用邪恶诡计，以致统治者从农夫或圣职人员手中购得土地，此类交易亦应无效，土地应归还主教区或隐修院。如此行的主教或院长应被驱逐，主教离开其主教区，院长离开其隐修院，因为他们浪费了并未积蓄之物。\par

\Emph{古代摘要：依神圣宗徒所赞同的，主教或隐修院院长对教区或隐修院财物所作的任何转让行为无效。如此行的主教或院长应被驱逐。}\par

% structure: p0201 source=h3 target=subsection reason=relative-heading-rank; base=h2

\subsection{第十三条}

\Emph{那些将隐修院变为公共房屋者应受特别谴责。}\par

因我们的罪恶，在教会中发生的灾难期间，一些圣所，例如主教府和隐修院，被某些人夺取，变成了公共客栈。如果现在占有它们的人愿意归还，使其恢复原来用途，善莫大焉；否则，若这些人在圣职名单上，我们命令予以撤职；若他们是隐修士或平信徒，则处以绝罚，作为被圣父、圣子、圣神所定罪的人，并被安置在\BibleQuote{虫子不死，火也不灭}之处，因为他们违抗了主的声音，主说：\BibleQuote{不要使我父的殿成为商场}。\par

\Emph{古代摘要：凡擅用教区或隐修院公产者，若不将属于教区或隐修院之物归还主教或院长，则整个行为无效。若他们是圣职人员，应予以撤职；若为平信徒或隐修士，则予以绝罚。}\par

% structure: p0205 source=h3 target=subsection reason=relative-heading-rank; base=h2

\subsection{第十四条法令}

\Emph{凡未经祝圣者，不得在礼仪聚集期间的\OriginalTerm{读经台}{ambo}上诵读。}\par

众所皆知，司铎职中有一定的秩序，且谨慎维护司铎的晋升，甚得天主喜悦。因此，我们见到某些青年虽已接受剪发礼，却尚未从主教手中接受祝圣，便在礼仪聚集期间的\OriginalTerm{读经台}{ambo}上诵读，且如此违反法令，我们禁止今后再有此举，此禁令亦适用于隐修士。准许每位\OriginalTerm{院长}{hegumenos}在自己的隐修院内任命一位诵读员，只要他本人曾由主教覆手得\OriginalTerm{院长}{hegumenos}品位，且已知是司铎。\OriginalTerm{乡村主教}{chorepiscopi}亦得依古例，经主教授权，任命诵读员。\par

\Emph{任何人若非由主教祝圣，不得在\OriginalTerm{读经台}{ambon}上诵读。此规定亦适用于隐修士。隐修院院长若由主教祝圣，可在其院内任命诵读员，但仅限于其隐修院。\OriginalTerm{乡村主教}{chorepiscopus}亦可任命诵读员。}\par

% structure: p0209 source=h3 target=subsection reason=relative-heading-rank; base=h2

\subsection{第十五条法令}

\Emph{神职人员不应被派往两个教会。}\par

今后，任何神职人员不得被任命管理两个教会，因为这带有买卖和可耻牟利的色彩，且完全违背教会传统。我们亲耳听到主的声音说：\BibleQuote{没有人能事奉两个主人，因为他要么恨这个爱那个，要么依附这个而轻视那个}。因此，如宗徒所说，各人应在蒙召时的身份中安分，并只在一所教会中服务。因为在教会事务中，通过可耻牟利所得，完全与天主隔绝。为满足生活所需，有各种职业，若有人愿意，可藉此谋取身体所需。因为宗徒说：\BibleQuote{这些手供应了我及同我一起之人的需要}。这类职业可在天主保护的城中寻得。但在城外乡间，因人数稀少，可予以宽免。\par

\Emph{古代摘要：今后在\Place{君士坦丁堡}，神职人员不得服务两个教会。但在郊外，因人数稀少，可予准许。}\par

% structure: p0213 source=h3 target=subsection reason=relative-heading-rank; base=h2

\subsection{第十六条法令}

\Emph{圣职人员不得穿着华丽服饰。}\par

一切轻浮举动和身体装饰，与司铎身份极不相称。因此，那些穿着艳丽华丽服饰的主教和神职人员，应当改正自己；若不改正，应受惩罚。同样，那些涂抹香水的人也当如此。当苦根滋生时，诽谤基督徒者的异端污染被注入天主教会。凡被此污染玷污者，不仅憎恶画像，而且藐视一切端庄，对生活严肃虔诚的人极其恼怒，以致应验了经上所写的：\Quote{罪人厌恶事奉天主}。因此，若有人被发现嘲笑那些穿着朴素庄重衣服的人，应以惩罚加以纠正。因为从古以来，每位圣职者都穿着朴素庄重的衣服；确实，凡衣着不仅因需要，而是为外表炫耀，便带有浮华之气，如\Person{大巴西略}所言。亦无人穿着绣花丝绸衣服，或在衣边上装饰彩色饰物；因为他们曾从主的口中听到：\BibleQuote{穿华丽衣服的人是在王宫里}。\par

\Emph{古代摘要：穿着华丽衣服、涂抹香水的主教和神职人员应受纠正。若执迷不悟，应受惩罚。}\par

% structure: p0217 source=h3 target=subsection reason=relative-heading-rank; base=h2

\subsection{第十七条法令}

\Emph{若无能力完成工程者，不得开始建造小堂。}\par

某些隐修士因贪图管辖而离开隐修院，又不愿服从，便着手建造小堂，但无财力完成。凡擅自进行此类工程者，应由当地主教予以禁止。但若其有财力完成所计划之事，则应让其继续完成。此规则亦适用于平信徒和神职人员。\par

\Emph{古代摘要：凡愿建造隐修院者，若有财力完成，应开始工作并使其完工；否则，应由当地主教禁止。此规则亦适用于平信徒和隐修士。}\par

% structure: p0221 source=h3 target=subsection reason=relative-heading-rank; base=h2

\subsection{第十八条法令}

\Emph{妇女不宜住在主教府或男子隐修院内。}\par

神圣宗徒说：\Quote{对外人不可有失检点。}\newline{}如今妇女住在主教住所或隐修院内，已为许多人制造了严重失足的机会。因此，凡明知在主教的府邸或隐修院内留用女奴或自由妇女以执行某些任务者，应受斥责。若仍继续留用，则应受撤职处分。如有妇女在主教的郊区庄园工作，主教或\OriginalTerm{院长}{hegumenos}欲前往该处时，在主教或院长在场期间，该妇女不得继续工作，而应退避到别处，直至主教（或院长）离开，以免引起抱怨。\par

\Emph{《古代节录》：妇女不应被留置在主教的住所或隐修院内。若有胆敢如此行事的，应受谴责；若坚持不改，应受撤职处分。凡主教或隐修院院长在场时，任何妇女不得侍奉或露面，而应退避，直至他离开。}\par

% structure: p0225 source=h3 target=subsection reason=relative-heading-rank; base=h2

\subsection{第十九条}

\Emph{领受圣秩者、隐修士及修女所发的誓愿，不得有任何财物的索取。}\par

可耻的贪财之恶已在教会主管者中如此猖獗，以至于某些所谓虔诚的男女，竟完全忘记了主的诫命，受了迷惑，为金钱而接纳那些自荐接受圣秩或隐修生活的人。因此，如此被接纳者的第一步既属不法，整个过程便成无效，正如大圣巴西略所言。因为天主不能藉\OriginalTerm{玛门}{mammon}来侍奉。所以，若有人发现做了这类事，无论他是主教、\OriginalTerm{院长}{hegumenos}或司铎，要么立即停止，否则应受撤职处分，按照神圣卡尔西顿大公会议第二条教规。若犯者是一位\OriginalTerm{女院长}{abbess}，应把她从她的隐修院迁出，安置到另一隐修院中，居于从属地位。同样，对一位未受司铎圣职的院长，也应如此处理。至于父母为子女作为嫁妆所捐赠的，或本人从自己财产中奉献的，并声明是献给天主的礼物，我们规定：不论受益者是否继续留在隐修院，这些礼物应按照最初的决定，归隐修院所有；除非确实有理由控告其上司。\par

\Emph{《古代节录》：凡为金钱而接纳前来领受圣秩或隐修生活者，若他是主教、隐修院院长或其他圣职人员，要么停止，要么撤职。隐修院的女院长（若行此事）应被驱逐，并置于服从之下。隐修士的院长若非司铎，亦同。但来者所带来的财产，无论其人是否留下，应保持不变，除非院长有过错。}\par

% structure: p0229 source=h3 target=subsection reason=relative-heading-rank; base=h2

\subsection{第二十条}

\Emph{此后不得再建立男女混合的隐修院；并论及已有的混合隐修院。}\par

我们规定：此后不得再建立男女混合的隐修院，因为这已成为许多人的诟病和抱怨的原因。对于那些愿与家人一起离世而追随隐修生活的人，让男子进入男隐修院，女子进入女隐修院，因为这是中悦天主的。现有的混合隐修院，应遵守我们圣父巴西略的规则，并依其训令管理。隐修士与修女不得同住一所隐修院，因同住一处会为奸淫制造机会。隐修士不得进入女隐修院；修女也不得进入男隐修院与任何人交谈。隐修士不得睡在女隐修院内，也不得单独与修女共餐。当男子为\OriginalTerm{咏经修女}{canonesses}送食物时，女院长应偕同一位年长修女在女院门外接收。当隐修士想要探望其在女隐修院中的女亲属时，应在女院长面前与她交谈，且言语简短，然后迅速离开。\par

\Emph{《古代节录》：不得有混合隐修院，隐修士与修女不得同住一建筑，也不得单独交谈。此外，若有人带东西给咏经修女，应在外等候并交给她；他应在女院长面前会见亲属。}\par

% structure: p0233 source=h3 target=subsection reason=relative-heading-rank; base=h2

\subsection{第二十一条}

\Emph{隐修士不得离开自己的隐修院而投奔其他隐修院。}\par

隐修士或修女不应离开所属的隐修院而投奔其他隐修院。若有人这样做，应把他当作客人接待。然而，未经其\OriginalTerm{院长}{hegumenos}同意，不得使他成为该院的成员。\par

\Emph{《古代节录》：隐修士或修女不得离开自己的隐修院而进入另一隐修院；若他（或她）进入，应作为客人接待；但不得违背院长的意愿被接纳或得到款待。}\par

% structure: p0237 source=h3 target=subsection reason=relative-heading-rank; base=h2

\subsection{第二十二条}

\Emph{当隐修士偶尔需与妇女共餐时，应遵守谢恩、节制和明智之道。}\par

将万事交托于天主，不随从自己的私意，实为大利。因为神圣的宗徒说：\BibleQuote{你们或吃或喝，一切都要为光荣天主而行。}\BibleRef{格前 10:31} 基督我们的天主也在福音中命令我们，要断绝罪的根源；因为他不仅谴责奸淫的行为，连心中趋向奸淫的动念也加以谴责，他说：\BibleQuote{凡注视妇女，有意贪恋她的，他已在心里奸淫了她。}\BibleRef{玛 5:28} 我们既受了这样的教训，就应当洁净自己的心思。固然一切事都是许可的，但并非一切事都有益处，这是我们从宗徒的口中受教的。众人必须饮食才能生存。对于那些以结婚生子及平信徒身份为生活的人，男女一同进食并无不当，只须感谢赐饮食者即可；但若有剧场式的娱乐，即伴以竖琴淫靡曲调的撒殚歌曲，那么先知所说的咒诅便会临到他们身上，他说：\BibleQuote{祸哉，那些弹琴奏乐饮酒，却不注视主的作为，不关心他手所行的人！}\BibleRef{依 5:11-12} 凡在基督徒中发现此类行为者，应自改正；若不肯改正，则前代教父在教规中所规定的处罚必临到他们。至于那些生活无忧、远离世俗的人，即那些在主天主面前决意背负独身轭的人，当独自静默而坐。此外，凡选择司铎生活的人，私下与妇女共餐是绝对不合法的，除非有敬畏天主且谨慎的男女在场，好使他们的宴席也能转化为灵性的建树。对亲戚亦应同样遵守。再者，若有隐修士或司铎在旅途中未带必备之物，因急需而认为可以投宿客栈或某人家中，这是许可的，因为需要使然。\par

\Emph{古代摘要：平信徒女子与男子共餐并无不妥；但选择独身生活的男子私下与妇女共餐则不合宜；除非与敬畏天主的人及虔诚的男女一同。但在旅行时，隐修士或神职人员若未携带必需用品，可进入客栈。}\par

% structure: p0241 source=h2 target=section reason=relative-heading-rank; base=h2

\section{大公会议致皇帝与皇后信函}

（拉贝与科萨尔，《会议集》，第七卷，第577栏。）\par

致我们至为虔诚与至为温和的君王，君士坦丁及其母后伊琳。塔拉西乌斯，你们受天主保护的皇城新罗马的不配主教，以及全体神圣会议——此会议因天主的美意并奉你们爱基督的陛下的命令，在著名的尼西亚都城召开，为在此城第二次集会——谨祝。\par

基督，我们的天主（他是教会的头），极尊贵的君王，当你们那握在他手中的心说出那句美言，命我们因他的名聚集时，他便受到了光荣，为使我们能坚定持守教会信理中那确凿、不变且由天主所赐的真理。你们的头戴金冠和极璀璨的宝石，同样，你们的心思也以福音的训诫和教父的教导为装饰。你们既为那些声音传遍普世者的门徒和同伴，就成了所有基督名下之人在虔诚道路上的领袖，清晰地陈述了真理之言，并树立了正统与虔诚的光辉榜样；因此你们在信友面前如同燃烧的明灯。那濒临倾覆的教会，你们亲手扶助，以纯正的教理坚固她，并使纷争的人们归于正确判断的合一。因此我们可以大胆地说，是天主的美意藉着你们使虔敬得胜，并使我们口中充满喜乐，舌上充满欢欣。这些话我们以正式法令从唇间说出。因为还有什么比维护教会的利益更光荣？还有什么更能激发我们的喜乐呢？\par

有些人兴起，具有虔诚的外形，因为他们披着司铎尊严的外衣，却否认其能力；因此他们为自己招来只是\Place{巴比伦}司铎的指责。关于这样的人，预言的话早已宣告说：\Quote{不法之事从\Place{巴比伦}的司铎中发出。}不仅如此，他们还聚集在一起，形成一个如同\Person{盖法}所主持的公议会，成为不敬虔教义的传播者。他们满口咒骂和苦毒，以为用辱骂的话语就能赢得上风。用毁谤的舌头和同样性质的笔，反对天主亲自使用的术语，他们编造了奇妙的故事，然后着手将王室司祭团和圣洁的国民——即那些穿上基督、并藉他的恩宠得以避免偶像愚妄的人——污蔑为偶像崇拜者。他们心存恶念，从事不法之事，企图完全压制可敬圣像的描绘。因此，凡是镶嵌画制成的圣像，他们挖掉；凡是蜡彩画的，他们刮去；从而使圣殿的优美完全陷入混乱。在这些行为中，特别值得注意的是，他们将那些为纪念我们的天主基督和他的圣人而设立在板上的画像投入火中。总之，一言以蔽之，他们亵渎了我们的教堂，使之陷入彻底混乱。然后，有些主教成为这一异端的首领，在原本和平的地方，他们在人民中煽动纷争；在教会的田地里，他们不播种麦子，反而播种莠子。他们将酒与水混合，将污秽的饮料递给周围的人。他们虽然是\Place{阿拉伯}的豺狼，却披着羊皮藏匿，用似是而非的推理反对真理，试图推销他们的谎言。但先知说，他们始终\BibleQuote{孵育毒蛇的蛋，编织蜘蛛的网}；\BibleQuote{谁要吃他们的蛋，打碎一个，发现是坏蛋，里面有毒蛇，}散发出致命的臭气。\par

在这样的情况下，谎言正在毁灭真理，你们，最仁慈、最高贵的君王，并没有无所作为地允许如此严重的灾祸和如此毁灭灵魂的错误在你们的时代长久延续。而是被住在你们内的天主圣神所感动，你们竭尽全力彻底铲除它，从而保持教会管理的稳定，以及臣民之间的和谐；使你们的整个帝国得以在和平中建立，与你们的名号相符。你们正确地判断，不能容忍的是，在其他事情上我们能同心合意、和睦相处，但在我们生命中本应首要关心的事——教会的和平——上，我们之间却有纷争和分裂。而当我们以基督为元首时，我们本应藉着彼此的合一和信德，互为肢体，成为一个身体。因此，你们命令我们神圣而人数众多的会议在\Place{尼凯亚}大都市聚集，以便在清除教会的分裂之后，我们将分离的肢体重新合为一体，并小心撕碎并彻底毁灭他们用荆棘纤维编织的谬论外衣，重新展开正统信仰的美丽袍服。\par

如今，我们仔细追溯宗徒和教父的传统，敢于发言。藉着至圣圣神的嘘气，我们同心合意，完全合一，理解天主教会（原文\Quote{大公教会}应改为\Quote{天主教会}）的和谐传统，我们与六次神圣大公会议所提出的谐音完全一致；因此，我们诅咒了\Person{亚略}的疯狂、\Person{马其顿纽}的狂热、\Person{阿波利纳里}的无理性理解、\Person{聂斯多略}的崇拜人性、\Person{欧提基}和\Person{狄约斯哥罗}所发明的对于本性的不敬混合，以及与他们同行的多头水蛇。我们也诅咒了\Person{奥利振}、\Person{狄第默}和\Person{厄瓦格留}的虚妄故事；以及\Person{塞尔吉}、\Person{何诺留}、\Person{西禄}和\Person{彼禄}所主张的一志论，更确切地说，我们诅咒了他们自己的邪恶意志。最后，受圣神教导，从他那里我们汲取了清水，我们同心合意，以同一灵魂，完全用神圣教义的海绵擦去了新近发明的异端，这异端堪与刚才提到的那些并列，它在它们之后兴起，关于神圣圣像发出如此空洞的胡说。我们将这些虚妄却具有颠覆性的空谈的制造者远远驱逐出教会的地界。\par

就如手足依心智的吩咐而动，同样，我们既领受了圣神的恩宠与力量，又得着你们皇权的协助与合作，便异口同声地宣认并宣告为真理：我们主耶稣基督的圣像——因祂确是真人——以及描绘福音所历史记述的那些圣像，还有我们无玷的圣母、天主之母的圣像，同样，圣天使的圣像（因为他们曾向那些堪当看见异象的人显现于人形），以及任何圣人的圣像，都应当保存并保留。\newline{}【我们还规定】圣人的英勇事迹应绘于板画上、墙壁上、圣器与祭衣上，正如从天主圣教会自古以来的习惯；这习惯在教会初期诸圣领袖的训导中，以及他们之后的可敬教父们的训导中，被视为具有法律效力。\newline{}【我们也同样规定】这些图像应受\OriginalTerm{敬礼}{\GreekText{προσκυνεῖν}}，即向它们致以问候。使用这个词语的理由是，它具有双重含义。因为在古希腊语中，\OriginalTerm{问候或亲吻}{\GreekText{κυνεῖν}}既表示\Quote{问候}，也表示\Quote{亲吻}。而介词\OriginalTerm{朝向}{\GreekText{προς}}则给它增添了指向对象的热切渴望之意；例如，我们有\OriginalTerm{携带}{\GreekText{φέρω}}和\OriginalTerm{呈上}{\GreekText{προσφέρω}}，\OriginalTerm{相遇}{\GreekText{κυρῶ}}和\OriginalTerm{前来相遇}{\GreekText{προσκυρῶ}}，同样也有\OriginalTerm{问候}{\GreekText{κυνέω}}和\OriginalTerm{敬礼}{\GreekText{προσκυνέω}}。这后一个词意味着问候与热切之爱；因为一个人所爱的，他也\OriginalTerm{敬礼}{\GreekText{προσκυνεῖ}}，而他所敬礼的，他也极其爱慕，正如我们日常对待所爱之人的习惯所证实的那样，当两位朋友相遇时，这两种含义都在实践中体现。这个词不仅为我们所用，我们也在古人的神圣圣经中见到它被记载。因为列王纪上写道：\Quote{达味就起身，面伏于地，向约纳堂\OriginalTerm{敬礼}{\GreekText{προσεκυνήσε}}三次，并亲吻了他}\BibleRef{列上 20:41}。主在福音中论及法利塞人说了什么？\Quote{他们喜爱筵席上的首位，街市上的\OriginalTerm{问候}{\GreekText{ἀσπασμοὺς}}}。显然，这里的\Quote{问候}是指\OriginalTerm{敬礼}{\GreekText{προσκύνησιν}}，因为法利塞人极其高傲，自以为是义人，渴望被所有人敬礼，而不仅仅是亲吻。因为接受后一种问候显得过于谦卑，不合法利塞人的心意。我们还有神圣宗徒保禄的例证，路加在宗徒大事录中记载：\Quote{我们到了耶路撒冷，弟兄们欢欢喜喜地接待了我们。第二天，保禄同我们去见雅各伯，所有的长老也都在场。保禄\OriginalTerm{问候}{\GreekText{ἀσπασάμενος}}了他们，就一一述说天主藉他的职务在外邦人中所行的事}\BibleRef{宗 21:17}。这里所提到的问候，显然宗徒意在表示彼此间所行的\OriginalTerm{尊敬的敬礼}{\GreekText{τιμητικὴν} \GreekText{προσκύνησιν}}，正如他论及雅各伯时所说的：\Quote{他\OriginalTerm{敬礼}{\GreekText{προσεκύνησεν}}了自己的杖头}\BibleRef{希 11:21}。与这些例子相符的是，被称为神学家的额我略所说的：\Quote{尊敬白冷，并\OriginalTerm{敬礼}{\GreekText{προσκυνήσον}}马槽。}\par

如今，那些正确而真诚地理解圣经的人中，有谁曾以为我们所引用的这些例子指的是\OriginalTerm{在圣神内的崇拜}{\GreekText{τῆς} \GreekText{ἐν} \GreekText{πνεύματι} \GreekText{λατρείας}}呢？【肯定无人如此想过，】除非有些人全然失去理智，对圣经和教父训导一无所知。雅各伯岂会\OriginalTerm{崇拜}{\GreekText{ἐλάτρευσεν}}他的杖头？神学家额我略岂会命令我们\OriginalTerm{崇拜}{\GreekText{λατρεύειν}}马槽？决不是。\par

再者，当我们向赋予生命的十字架致以问候时，我们一同歌唱：\Quote{我们\OriginalTerm{敬礼}{\GreekText{προσκυνῶμεν}}你的十字架，主，我们也\OriginalTerm{敬礼}{\GreekText{προσκυνῶμεν}}那刺开你良善生命之侧的枪。}这显然只是一种问候，如此称呼，其实质在于我们用嘴唇接触所提及之物。我们承认，\OriginalTerm{敬礼}{\GreekText{προσκύνησις}}一词在圣经和我们博学圣洁的教父著作中，频繁地用于\OriginalTerm{在圣神内的崇拜}{\GreekText{ἐπὶ} \GreekText{της} \GreekText{ἐν} \GreekText{πνεύματι} \GreekText{λατρείας}}，因为作为一个多义词，它可以用来表达那种属于服役的尊敬。此外还有出于尊敬、爱慕和敬畏的敬礼。正是在这个意义上，我们敬拜你们光荣而最尊贵的威严。同样，还有另一种仅出于害怕的敬礼，如雅各伯敬礼厄撒乌。还有感恩的敬礼，如亚巴郎向赫特人致敬，因他为他们所赐的埋葬妻子撒辣的田地。最后，那些期望获得恩赐的人，向高于他们的人致敬，如雅各伯敬礼法郎。因此，由于这个词有这么多含义，圣经教训我们说：\Quote{你当敬礼主你的天主，唯独事奉祂}，它只说敬礼应给予天主，但并未加\Quote{唯独}一词，因为敬礼作为含义广泛的词，是一个含糊的术语；但它接着说：\Quote{你当\OriginalTerm{事奉}{\GreekText{λατρεύσεις}}祂独一}，因为我们只向天主行\OriginalTerm{钦崇}{latria}。\par

我们所规定的这些事既如此有力地得到支持，那么主\Person{耶稣基督}作为人的圣像、以及无玷的天主之母终身童贞\Person{玛利亚}的圣像、可敬的天使和众圣人的圣像应受敬礼和问候，这在天主面前显然是可接受的并中悦天主的。若有人不信如此，反而妄图进一步争论，并对圣像应有的敬礼心怀恶意，那么我们的神圣大公会议（藉天主圣神的内在运作，以及教父与教会的传统所坚固）就予以绝罚。绝罚无非是完全与天主分离。因为若有人好争论，不肯顺服接受现在所规定的事，他们就是踢刺棒，并在与\Person{基督}作战中伤害自己的灵魂。他们以加给教会的侮辱为乐，显然表明自己正是那些疯狂攻击虔敬的人，因此应被视为与古时的异端者同类，以及他们不敬神的同伴和弟兄。\par

我们派遣了我们可敬的弟兄和同工司铎，即蒙天主所爱的主教们，以及一些修道院长和神职人员，以便他们将我们的议事详尽报告给您的虔敬耳中。为证实和确认我们所规定的，并为您最虔敬的陛下提供保证，我们呈上了来自教父的证据，我们收集了许多证据以阐明这灿烂的真理，这里仅呈上少许。\par

愿我们众人的救主，就是\OriginalTerm{与你们一同统治}{\GreekText{συμβασιλεύων} \GreekText{ὑμῖν}}、且愿意藉你们将祂的平安赐予教会的那一位，保全你们的王国多年，也保全你们的议会、王公、忠信的军队，以及帝国的整个国度；并愿祂赐给你们战胜所有仇敌的胜利。因为正是祂说：\BibleQuote{我指着我的永生起誓，上主说：那光荣我的，我必光荣。}也是祂赐给你们力量，将打击你们的一切仇敌，并使你们的子民欢乐。\par

哦，新城\Place{熙雍}，你当欢喜踊跃；你是普世的奇观。因为虽然\Person{达味}没有在你内为王，然而你虔敬的君王在此治理你的事务，如同\Person{达味}所做的一样。\Person{主}在你中间；愿祂的名永受赞美，世世无穷。阿们。\par

\SourceRecord{Translated by Henry Percival.；Nicene and Post-Nicene Fathers, Second Series；Vol. 14.；Edited by Philip Schaff and Henry Wace.；Buffalo, NY: Christian Literature Publishing Co.,；1900.；<http://www.newadvent.org/fathers/3819.htm>.}
